% ─────────────────────────────────────────────────────────────────────────────
%  Cadenas de Markov Homogéneas — Modelos Estocásticos
%  Autor: P. Hurtado · 2026
% ─────────────────────────────────────────────────────────────────────────────
\documentclass[10pt,a4paper]{article}

\usepackage[T1]{fontenc}
\usepackage[utf8]{inputenc}
\usepackage{lmodern}
\usepackage[a4paper, top=1.8cm, bottom=2.0cm, left=1.5cm, right=1.5cm]{geometry}
\usepackage[spanish,es-noshorthands]{babel}
\usepackage{amsmath,amssymb,mathtools,amsthm}
\usepackage{xcolor}
\usepackage[most]{tcolorbox}
\usepackage{tikz}
\usepackage{pgfplots}
\pgfplotsset{compat=1.18}
\usepackage{microtype}
\usepackage{needspace}
\usepackage{parskip}
\usepackage{fancyhdr}
\usepackage{hyperref}
\usepackage{booktabs}
\usepackage{array}
\usepackage{caption}

\usetikzlibrary{arrows.meta, positioning, calc, shapes.geometric, fit, backgrounds, decorations.pathreplacing}

\setlength{\parskip}{3pt}
\setlength{\parindent}{0pt}
\setlength{\headheight}{19pt}

% ─── Numeración de ecuaciones por sección ────────────────────────────────────
\numberwithin{equation}{section}

% ─── Colores ─────────────────────────────────────────────────────────────────
\definecolor{sectionbg}{rgb}{0.16,0.26,0.52}
\definecolor{subsecbg}{rgb}{0.38,0.50,0.75}
\definecolor{codeblue}{rgb}{0.0,0.28,0.68}
\definecolor{notebg}{rgb}{1.0,0.97,0.87}
\definecolor{noteborder}{rgb}{0.80,0.65,0.25}
\definecolor{defbg}{rgb}{0.93,0.98,0.93}
\definecolor{defborder}{rgb}{0.25,0.60,0.30}
\definecolor{thmbg}{rgb}{0.92,0.95,1.0}
\definecolor{thmborder}{rgb}{0.25,0.45,0.80}
\definecolor{exbg}{rgb}{0.96,0.96,0.96}
\definecolor{exborder}{rgb}{0.55,0.55,0.55}
\definecolor{warnbg}{rgb}{1.0,0.94,0.93}
\definecolor{warnborder}{rgb}{0.80,0.25,0.20}
\definecolor{proofbg}{rgb}{0.98,0.97,1.0}
\definecolor{proofborder}{rgb}{0.55,0.40,0.75}

% ─── Hyperref ────────────────────────────────────────────────────────────────
\hypersetup{
  colorlinks=true,
  linkcolor=codeblue,
  citecolor=codeblue,
  urlcolor=codeblue,
}

% ─── Cabeceras ───────────────────────────────────────────────────────────────
\pagestyle{fancy}
\fancyhf{}
\fancyhead[L]{\footnotesize\sffamily\color{sectionbg}Resumen de Probabilidad --- Modelos Estoc\'asticos}
\fancyhead[R]{\footnotesize\sffamily\color{sectionbg}2026}
\fancyfoot[C]{\footnotesize\sffamily\thepage}
\renewcommand{\headrulewidth}{0.4pt}
\renewcommand{\headrule}{\hbox to\headwidth{\color{subsecbg}\leaders\hrule height \headrulewidth\hfill}}

% ─── tcolorbox styles ────────────────────────────────────────────────────────
\tcbset{
  defbox/.style={
    enhanced,
    colback=defbg, colframe=defborder,
    boxrule=0.6pt, arc=3pt,
    left=7pt, right=7pt, top=5pt, bottom=5pt,
    fonttitle=\small\bfseries\sffamily,
    coltitle=white, colbacktitle=defborder,
    titlerule=0pt,
    colupper=black,
  },
  thmbox/.style={
    enhanced,
    colback=thmbg, colframe=thmborder,
    boxrule=0.6pt, arc=3pt,
    left=7pt, right=7pt, top=5pt, bottom=5pt,
    fonttitle=\small\bfseries\sffamily,
    coltitle=white, colbacktitle=thmborder,
    titlerule=0pt,
  },
  exbox/.style={
    enhanced,
    colback=exbg, colframe=exborder,
    boxrule=0.4pt, arc=3pt,
    left=7pt, right=7pt, top=5pt, bottom=5pt,
    fonttitle=\small\bfseries\sffamily,
    coltitle=white, colbacktitle=exborder,
    titlerule=0pt,
  },
  notebox/.style={
    colback=notebg, colframe=noteborder,
    boxrule=0.4pt, arc=3pt,
    left=5pt, right=5pt, top=4pt, bottom=4pt,
    fontupper=\footnotesize\itshape,
  },
  warnbox/.style={
    enhanced,
    colback=warnbg, colframe=warnborder,
    boxrule=0.6pt, arc=3pt,
    left=7pt, right=7pt, top=5pt, bottom=5pt,
    fonttitle=\small\bfseries\sffamily,
    coltitle=white, colbacktitle=warnborder,
    titlerule=0pt,
    title={$\triangle$\ Atenci\'on},
    fontupper=\small,
  },
  proofbox/.style={
    enhanced,
    colback=proofbg, colframe=proofborder!60,
    boxrule=0.6pt, arc=3pt,
    left=7pt, right=7pt, top=5pt, bottom=5pt,
    fonttitle=\small\itshape\sffamily,
    coltitle=white, colbacktitle=proofborder!80,
    titlerule=0pt,
    title={Demostraci\'on},
    fontupper=\small,
  },
}

% ─── Helpers ─────────────────────────────────────────────────────────────────
\newcommand{\secheader}[2]{%
  \refstepcounter{section}%
  \addcontentsline{toc}{section}{\protect\numberline{\thesection}#2}%
  \vspace{10pt}\needspace{6cm}\noindent%
  \begin{tcolorbox}[
    colback=sectionbg, colupper=white, colframe=sectionbg,
    sharp corners, left=8pt, right=8pt, top=5pt, bottom=5pt,
    fontupper=\large\bfseries\sffamily,
  ]
  Secci\'on #1:\enspace #2
  \end{tcolorbox}\vspace{2pt}%
}

\newcommand{\subsecheader}[1]{%
  \needspace{4cm}\vspace{6pt}\noindent%
  \addcontentsline{toc}{subsection}{#1}%
  \begin{tcolorbox}[
    colback=subsecbg!18, colframe=subsecbg!45,
    sharp corners, boxrule=0.4pt,
    left=5pt, right=5pt, top=2pt, bottom=2pt,
    fontupper=\small\bfseries\sffamily,
  ]
  #1
  \end{tcolorbox}\vspace{1pt}%
}

% ─── Macros matemáticos (módulo 01) ──────────────────────────────────────────
\newcommand{\Prob}[1]{\mathbb{P}\!\left(#1\right)}
\newcommand{\E}[1]{\mathbb{E}\!\left[#1\right]}
\newcommand{\Var}[1]{\mathrm{Var}\!\left(#1\right)}
\newcommand{\Cov}[2]{\mathrm{Cov}\!\left(#1,\,#2\right)}
\newcommand{\sigalg}{\mathcal{F}}
\newcommand{\ind}{\mathbf{1}}
\newcommand{\Borel}{\mathcal{B}}
\newcommand{\R}{\mathbb{R}}
\newcommand{\Nset}{\mathbb{N}}
\DeclareMathOperator{\supp}{supp}

% ─── Macros adicionales para cadenas de Markov ───────────────────────────────
\newcommand{\Sset}{\mathcal{S}}                          % espacio de estados
\newcommand{\Pmat}{\mathbf{P}}                           % matriz de transición
\newcommand{\generator}{\mathbf{Q}}                      % generador infinitesimal
\newcommand{\pij}[2]{p_{\,#1#2}}                         % prob. de transición 1 paso
\newcommand{\pijn}[3]{p_{\,#1#2}^{(#3)}}                 % prob. en n pasos
\newcommand{\fij}[2]{f_{\,#1#2}}                         % prob. de primer paso
\newcommand{\muu}[1]{\mu_{#1}}                           % tiempo medio de retorno

% ─── Contadores ──────────────────────────────────────────────────────────────
\newcounter{defcount}[section]
\renewcommand{\thedefcount}{\thesection.\arabic{defcount}}
\newcounter{thmcount}[section]
\renewcommand{\thethmcount}{\thesection.\arabic{thmcount}}
\newcounter{excount}[section]
\renewcommand{\theexcount}{\thesection.\arabic{excount}}

% ─── Entornos ────────────────────────────────────────────────────────────────
\newenvironment{definicion}[1]{%
  \refstepcounter{defcount}%
  \begin{tcolorbox}[defbox, title={Definici\'on \thedefcount\ \textperiodcentered\ #1}]%
}{\end{tcolorbox}}

\newenvironment{teorema}[1]{%
  \refstepcounter{thmcount}%
  \begin{tcolorbox}[thmbox, title={Teorema \thethmcount\ \textperiodcentered\ #1}]%
}{\end{tcolorbox}}

\newenvironment{proposicion}[1]{%
  \refstepcounter{thmcount}%
  \begin{tcolorbox}[thmbox, title={Proposici\'on \thethmcount\ \textperiodcentered\ #1}]%
}{\end{tcolorbox}}

\newenvironment{ejemplo}[1]{%
  \refstepcounter{excount}%
  \begin{tcolorbox}[exbox, title={Ejemplo \theexcount\ \textperiodcentered\ #1}]%
}{\end{tcolorbox}}

\newenvironment{nota}{\begin{tcolorbox}[notebox]}{\end{tcolorbox}}

\newenvironment{cuidado}{\begin{tcolorbox}[warnbox]}{\end{tcolorbox}}

\newenvironment{demostracion}{\begin{tcolorbox}[proofbox]}{\hfill$\square$\end{tcolorbox}}

% ─────────────────────────────────────────────────────────────────────────────
\begin{document}
% ─────────────────────────────────────────────────────────────────────────────

%  T I T U L O
\begin{tcolorbox}[
    colback=sectionbg, colupper=white, colframe=sectionbg,
    sharp corners, left=12pt, right=12pt, top=10pt, bottom=10pt,
  ]
  {\Large\bfseries\sffamily Cadenas de Markov Homog\'eneas}\\[4pt]
  {\normalsize\sffamily M\'odulo 2 --- Modelos Estoc\'asticos}\\[6pt]
  {\small\sffamily\color{white!80!sectionbg}
    Propiedad de Markov \textperiodcentered\ Tiempo discreto \textperiodcentered\
    Clasificaci\'on de estados \textperiodcentered\ Distribuciones estacionarias
    \textperiodcentered\ Tiempos de primer paso \textperiodcentered\
    Tiempo continuo \textperiodcentered\ Nacimiento y muerte}\\[6pt]
  {\footnotesize\sffamily\color{white!85!sectionbg}\itshape
    Por P. Hurtado \textperiodcentered\ 2026}
\end{tcolorbox}

\vspace{4pt}
\begin{nota}
  Este m\'odulo contin\'ua desde el M\'odulo~1 (probabilidad, variables aleatorias, proceso de Poisson).
  El Proceso de Poisson aparece como caso particular de cadena de Markov en tiempo continuo
  (proceso de nacimiento puro). Se asume familiaridad con \'algebra lineal (matrices, autovalores)
  y con la distribuci\'on Exponencial. Notaci\'on: $\Sset$ espacio de estados, $\Pmat$ matriz de
  transici\'on, $\pi$ distribuci\'on estacionaria, $\generator$ generador infinitesimal.
\end{nota}

\tableofcontents

\vspace{8pt}
\begin{tcolorbox}[
    colback=white, colframe=sectionbg,
    boxrule=0.6pt, arc=3pt,
    left=8pt, right=8pt, top=6pt, bottom=6pt,
    fonttitle=\small\bfseries\sffamily,
    title={Leyenda de colores},
  ]
  \small
  Cada bloque de contenido se identifica visualmente por color, seg\'un su naturaleza:
  \vspace{4pt}

  \renewcommand{\arraystretch}{1.35}
  \begin{tabular}{@{}p{0.7cm} l p{10.5cm}@{}}
    \tikz\fill[defbg, draw=defborder, line width=0.6pt] (0,0) rectangle (0.55,0.35);
                                          & \textbf{Definici\'on}      & Concepto formal con nombre y notaci\'on. Marco verde claro.                            \\
    \tikz\fill[thmbg, draw=thmborder, line width=0.6pt] (0,0) rectangle (0.55,0.35);
                                          & \textbf{Teorema/Proposici\'on} & Resultado demostrable a partir de las definiciones. Marco azul.                    \\
    \tikz\fill[proofbg, draw=proofborder!50, line width=0.6pt] (0,0) rectangle (0.55,0.35);
                                          & \textbf{Demostraci\'on}    & Argumento que justifica un teorema o proposici\'on. Marco morado.                      \\
    \tikz\fill[exbg, draw=exborder, line width=0.6pt] (0,0) rectangle (0.55,0.35);
                                          & \textbf{Ejemplo}           & Aplicaci\'on concreta o caso ilustrativo. Marco gris.                                  \\
    \tikz\fill[notebg, draw=noteborder, line width=0.6pt] (0,0) rectangle (0.55,0.35);
                                          & \textbf{Nota}              & Observaci\'on, intuici\'on o comentario complementario. Marco \'ambar.                 \\
    \tikz\fill[warnbg, draw=warnborder, line width=0.6pt] (0,0) rectangle (0.55,0.35);
                                          & \textbf{Atenci\'on}        & Advertencia sobre un error com\'un o sutileza. Marco rojo.                             \\
    \tikz\fill[sectionbg] (0,0) rectangle (0.55,0.35);
                                          & \textbf{Secci\'on}         & Encabezado principal del m\'odulo. Fondo azul oscuro.                                  \\
    \tikz\fill[subsecbg!18, draw=subsecbg!45, line width=0.6pt] (0,0) rectangle (0.55,0.35);
                                          & \textbf{Subsecci\'on}      & Subdivisi\'on tem\'atica dentro de una secci\'on. Fondo azul claro.                    \\
  \end{tabular}
\end{tcolorbox}

\newpage

% =============================================================================
\secheader{1}{Procesos Estoc\'asticos y Propiedad de Markov}

\subsecheader{1.1 \textperiodcentered\ Proceso estoc\'astico}

\begin{definicion}{Proceso estoc\'astico}
  Sea $(\Omega, \sigalg, P)$ un espacio de probabilidad y $T$ un conjunto de \'indices
  (conjunto de \emph{tiempo}). Un \textbf{proceso estoc\'astico} es una familia de variables
  aleatorias
  \[
    \{X_t : t \in T\},
  \]
  donde cada $X_t : \Omega \to \Sset$ toma valores en el \textbf{espacio de estados} $\Sset$.\\[3pt]
  Seg\'un la naturaleza de $T$ y $\Sset$ se distinguen cuatro tipos:
  \begin{center}
    \renewcommand{\arraystretch}{1.2}
    \begin{tabular}{@{}llll@{}}
      \toprule
      \textbf{Tiempo} $T$ & \textbf{Estados} $\Sset$ & \textbf{Tipo} \\
      \midrule
      Discreto ($\Nset$) & Discreto & Cadena en tiempo discreto \\
      Discreto ($\Nset$) & Continuo ($\R$) & Proceso en tiempo discreto \\
      Continuo ($[0,\infty)$) & Discreto & Cadena en tiempo continuo \\
      Continuo ($[0,\infty)$) & Continuo ($\R$) & Proceso en tiempo continuo \\
      \bottomrule
    \end{tabular}
  \end{center}
  En este m\'odulo $T = \Nset_0$ (discreto) o $T = [0,\infty)$ (continuo), y $\Sset$ es
  \emph{discreto} (finito o numerable).
\end{definicion}

\begin{nota}
  Para cada $\omega \in \Omega$ fijo, la funci\'on $t \mapsto X_t(\omega)$ es una
  \textbf{trayectoria} (o \emph{realizaci\'on}) del proceso.
  Un proceso estoc\'astico puede verse como una distribuci\'on sobre el espacio de trayectorias.
\end{nota}

\begin{center}
\begin{tikzpicture}[scale=1, >=Latex]
  \draw[->] (0,0) -- (8.5,0) node[right, font=\small] {$n$};
  \draw[->] (0,0) -- (0,3.5) node[above, font=\small] {$X_n$};
  \foreach \x in {0,1,2,3,4,5,6,7,8}
    \draw (\x, 0.05) -- (\x, -0.05) node[below, font=\tiny] {$\x$};
  \foreach \y/\lab in {1/$1$, 2/$2$, 3/$3$}
    \draw (0.05,\y) -- (-0.05,\y) node[left, font=\tiny] {\lab};
  % trayectoria ejemplo
  \draw[thick, color=sectionbg]
    (0,2) -- (1,2) -- (1,1) -- (2,1) -- (2,3) -- (3,3)
    -- (3,2) -- (4,2) -- (4,1) -- (5,1) -- (5,2) -- (6,2) -- (6,3)
    -- (7,3) -- (7,2) -- (8,2);
  \foreach \x/\y in {0/2,1/1,2/3,3/2,4/1,5/2,6/3,7/2,8/2}
    \fill[color=sectionbg] (\x,\y) circle (2.5pt);
  \node[font=\small\itshape, color=sectionbg!70!black] at (4,-0.8)
    {Trayectoria de una cadena en $\Sset = \{1,2,3\}$};
\end{tikzpicture}
\end{center}

\subsecheader{1.2 \textperiodcentered\ Filtraci\'on y adaptabilidad}

\begin{definicion}{Filtraci\'on y proceso adaptado}
  Una \textbf{filtraci\'on} $\{\sigalg_n\}_{n\ge0}$ es una sucesi\'on creciente de
  $\sigma$-\'algebras: $\sigalg_0 \subseteq \sigalg_1 \subseteq \cdots \subseteq \sigalg$.\\[3pt]
  Un proceso $\{X_n\}$ es \textbf{adaptado} a $\{\sigalg_n\}$ si $X_n$ es
  $\sigalg_n$-medible para todo $n$. Intuitivamente, $\sigalg_n$ representa la
  \emph{informaci\'on disponible hasta el instante $n$}.\\[3pt]
  La filtraci\'on natural del proceso es $\sigalg_n = \sigma(X_0, X_1, \ldots, X_n)$.
\end{definicion}

\subsecheader{1.3 \textperiodcentered\ Propiedad de Markov}

\begin{definicion}{Propiedad de Markov (d\'ebil)}
  Un proceso estoc\'astico $\{X_n\}_{n\ge0}$ con espacio de estados discreto $\Sset$ tiene la
  \textbf{propiedad de Markov} (d\'ebil) si para todo $n \ge 0$ y todos los estados
  $i_0, i_1, \ldots, i_{n-1}, i, j \in \Sset$:
  \[
    \Prob{X_{n+1} = j \mid X_n = i,\, X_{n-1} = i_{n-1}, \ldots, X_0 = i_0}
    = \Prob{X_{n+1} = j \mid X_n = i}.
  \]
  En palabras: el futuro depende del pasado \emph{\'unicamente} a trav\'es del presente.
  El estado actual $X_n$ captura toda la informaci\'on relevante del historial.
\end{definicion}

\begin{nota}
  La propiedad de Markov es una hip\'otesis de \emph{falta de memoria} del proceso.
  No dice que el futuro sea independiente del pasado en absoluto, sino que
  condicionar al presente hace que el pasado sea irrelevante.
\end{nota}

\begin{definicion}{Propiedad de Markov fuerte}
  Sea $\tau$ un \textbf{tiempo de parada} (variable aleatoria con valores en $\Nset_0$
  tal que $\{\tau = n\} \in \sigalg_n$). La \textbf{propiedad de Markov fuerte} establece que,
  dado $\tau < \infty$, el proceso desplazado $\{X_{\tau+k}\}_{k\ge0}$ tiene la misma
  distribuci\'on que $\{X_k\}_{k\ge0}$ iniciado en $X_\tau$, e independiente de
  $\sigalg_\tau$.\\[3pt]
  Toda cadena de Markov homog\'enea satisface la propiedad de Markov fuerte.
\end{definicion}

\subsecheader{1.4 \textperiodcentered\ Homogeneidad temporal}

\begin{definicion}{Cadena de Markov homog\'enea}
  Un proceso $\{X_n\}_{n\ge0}$ con la propiedad de Markov es \textbf{homog\'eneo en el tiempo}
  (o \emph{temporalmente homog\'eneo}) si las probabilidades de transici\'on no dependen de $n$:
  \[
    \Prob{X_{n+1} = j \mid X_n = i} = \pij{i}{j} \quad \forall\, n \ge 0,
  \]
  donde $\pij{i}{j}$ es la \textbf{probabilidad de transici\'on en un paso} de $i$ a $j$.
  En adelante toda cadena mencionada es homog\'enea salvo menci\'on expl\'icita.
\end{definicion}

\begin{cuidado}
  No confundir homogeneidad temporal con independencia. Los estados $X_n$ no son
  independientes; la homogeneidad solo garantiza que las \emph{probabilidades de salto} son
  constantes en el tiempo.
\end{cuidado}

% =============================================================================
\secheader{2}{Cadenas de Markov en Tiempo Discreto (CMTD)}

\subsecheader{2.1 \textperiodcentered\ Espacio de estados y probabilidades de transici\'on}

\begin{definicion}{Probabilidades de transici\'on en $n$ pasos}
  La \textbf{probabilidad de transici\'on en $n$ pasos} de $i$ a $j$ es
  \[
    \pijn{i}{j}{n} = \Prob{X_{n+m} = j \mid X_m = i}, \quad n \ge 1,
  \]
  que por homogeneidad no depende de $m$. Por convenio: $\pijn{i}{j}{0} = \delta_{ij}$
  (delta de Kronecker: 1 si $i=j$, 0 en caso contrario).
\end{definicion}

\subsecheader{2.2 \textperiodcentered\ Matriz de transici\'on}

\begin{definicion}{Matriz de transici\'on (matriz estoc\'astica)}
  Si $\Sset = \{1, 2, \ldots, s\}$ es finito, la \textbf{matriz de transici\'on} es la
  matriz $\Pmat \in \R^{s\times s}$ con entradas $(\Pmat)_{ij} = \pij{i}{j}$.\\[3pt]
  $\Pmat$ es una \textbf{matriz estoc\'astica por filas}: satisface
  \begin{enumerate}
    \item $\pij{i}{j} \ge 0$ para todo $i, j \in \Sset$.
    \item $\displaystyle\sum_{j \in \Sset} \pij{i}{j} = 1$ para todo $i \in \Sset$
          \quad (cada fila suma 1).
  \end{enumerate}
\end{definicion}

\begin{ejemplo}{Cadena del clima (2 estados)}
  Sea $\Sset = \{S, L\}$ (soleado, lluvioso). Si hoy es soleado, ma\~nana es soleado con
  probabilidad $0.8$ y lluvioso con $0.2$. Si hoy es lluvioso, ma\~nana es soleado con
  $0.4$ y lluvioso con $0.6$. La matriz de transici\'on es:
  \[
    \Pmat = \begin{pmatrix} 0.8 & 0.2 \\ 0.4 & 0.6 \end{pmatrix}
    \quad \leftarrow
    \begin{array}{l} \text{fila } S \\ \text{fila } L \end{array}
  \]
  Verificaci\'on: filas suman $1$. $\Pmat$ es estoc\'astica.
\end{ejemplo}

\begin{center}
\begin{tikzpicture}[>=Latex, node distance=3.5cm,
    state/.style={circle, draw=sectionbg, fill=thmbg, line width=1pt,
                  minimum size=1.1cm, font=\small\bfseries}]
  \node[state] (S) {$S$};
  \node[state, right of=S] (L) {$L$};
  \draw[->, thick, color=sectionbg]
    (S) edge[bend left=25] node[above, font=\small]{$0.2$} (L)
    (L) edge[bend left=25] node[below, font=\small]{$0.4$} (S);
  \draw[->, thick, color=defborder]
    (S) edge[loop left] node[left, font=\small]{$0.8$} (S)
    (L) edge[loop right] node[right, font=\small]{$0.6$} (L);
  \node[font=\small\itshape, color=sectionbg!70!black] at (2,-1.6)
    {Diagrama de transici\'on: cadena del clima};
\end{tikzpicture}
\end{center}

\subsecheader{2.3 \textperiodcentered\ Ecuaciones de Chapman-Kolmogorov}

\begin{teorema}{Ecuaciones de Chapman-Kolmogorov}
  Para todo $n, m \ge 0$ e $i, j \in \Sset$:
  \[
    \pijn{i}{j}{n+m} = \sum_{k \in \Sset} \pijn{i}{k}{n}\, \pijn{k}{j}{m}.
  \]
  En notaci\'on matricial: $\Pmat^{(n+m)} = \Pmat^{(n)} \Pmat^{(m)}$, con la convenci\'on
  $\Pmat^{(n)} = \Pmat^n$, donde $\Pmat^n$ es la $n$-\'esima potencia de $\Pmat$.
\end{teorema}

\begin{demostracion}
  Por la propiedad de Markov y marginalizaci\'on:
  \begin{align*}
    \pijn{i}{j}{n+m}
    &= \Prob{X_{n+m} = j \mid X_0 = i} \\
    &= \sum_{k \in \Sset} \Prob{X_{n+m}=j \mid X_n=k}\,\Prob{X_n = k \mid X_0 = i}
       \quad \text{(ley total)} \\
    &= \sum_{k \in \Sset} \pijn{k}{j}{m}\,\pijn{i}{k}{n}.
  \end{align*}
  Esto es exactamente la $(i,j)$-\'esima entrada del producto $\Pmat^n \Pmat^m = \Pmat^{n+m}$.
\end{demostracion}

\begin{nota}
  La clave es que \emph{la $n$-\'esima potencia de la matriz de transici\'on da las
  probabilidades en $n$ pasos}: $(\Pmat^n)_{ij} = \pijn{i}{j}{n}$.
\end{nota}

\subsecheader{2.4 \textperiodcentered\ Distribuci\'on en el paso $n$}

\begin{definicion}{Distribuci\'on marginal en el paso $n$}
  Sea $\boldsymbol{\mu}^{(0)}$ el vector fila de probabilidades iniciales:
  $\mu^{(0)}_i = \Prob{X_0 = i}$. La \textbf{distribuci\'on marginal} en el instante $n$ es
  \[
    \boldsymbol{\mu}^{(n)} = \boldsymbol{\mu}^{(0)}\,\Pmat^n,
  \]
  es decir, $\mu^{(n)}_j = \Prob{X_n = j} = \sum_{i\in\Sset} \mu^{(0)}_i\,\pijn{i}{j}{n}$.
\end{definicion}

\begin{ejemplo}{Distribuci\'on del clima en 2 pasos}
  Con la cadena del clima anterior y distribuci\'on inicial $(0.5, 0.5)$:
  \[
    \Pmat^2 = \begin{pmatrix}0.8&0.2\\0.4&0.6\end{pmatrix}^2
    = \begin{pmatrix}0.72&0.28\\0.56&0.44\end{pmatrix}.
  \]
  \[
    \boldsymbol{\mu}^{(2)} = (0.5,\,0.5)\begin{pmatrix}0.72&0.28\\0.56&0.44\end{pmatrix}
    = (0.64,\,0.36).
  \]
  Probabilidad de estar soleado a los 2 pasos: $0.64$.
\end{ejemplo}

\subsecheader{2.5 \textperiodcentered\ Forma can\'onica y clases de estados (anticipo)}

\begin{nota}
  Si el espacio de estados $\Sset$ se puede particionar en \emph{clases de comunicaci\'on}
  (ver Secci\'on~3), la matriz $\Pmat$ se puede reordenar en \emph{forma can\'onica}:
  \[
    \Pmat = \begin{pmatrix} \mathbf{Q} & \mathbf{R} \\ \mathbf{0} & \mathbf{I} \end{pmatrix},
  \]
  donde $\mathbf{Q}$ describe transiciones entre estados transitorios y $\mathbf{I}$
  los estados absorbentes. Esta forma es central para tiempos de absorci\'on (Secci\'on~5).
\end{nota}

\subsecheader{2.6 \textperiodcentered\ Propiedades estructurales de $\Pmat$}

\begin{proposicion}{Propiedades de $\Pmat$ y sus potencias}
  \begin{enumerate}
    \item \textbf{Estocasticidad hereditaria:} $\Pmat^n$ es estoc\'astica por filas $\forall\,n\ge0$. En particular $\Pmat^0=\mathbf{I}$.
    \item \textbf{Potencias = $n$-pasos:} $(\Pmat^n)_{ij}=\pijn{i}{j}{n}$.
    \item \textbf{Evoluci\'on de la distribuci\'on:} $\boldsymbol{\mu}^{(n)}=\boldsymbol{\mu}^{(0)}\Pmat^n$.
    \item \textbf{Autovalor trivial:} $\lambda=1$ es autovalor de $\Pmat$ con vector propio derecho $\mathbf{1}$. $\boldsymbol{\pi}$ estacionaria $\Leftrightarrow$ $\boldsymbol{\pi}\Pmat=\boldsymbol{\pi}$ (autovector izquierdo con $\lambda=1$).
    \item \textbf{Cota espectral:} $|\lambda|\le1$ para todo autovalor $\lambda$ de $\Pmat$.
  \end{enumerate}
\end{proposicion}

\begin{demostracion}
  \textbf{(1)} Inducci\'on sobre $n$. Base $n=0$: $\Pmat^0=\mathbf{I}$, claramente estoc\'astica.
  Paso inductivo: dado $\Pmat^{n-1}$ estoc\'astica,
  $(\Pmat^n)_{ij}=\sum_k(\Pmat^{n-1})_{ik}\pij{k}{j}\ge0$ y
  $\sum_j(\Pmat^n)_{ij}=\sum_k(\Pmat^{n-1})_{ik}\underbrace{\sum_j\pij{k}{j}}_{=1}=1$.\\[4pt]
  \textbf{(2)} Tambi\'en por inducci\'on. Para $n=1$: $(\Pmat)_{ij}=\pij{i}{j}=\pijn{i}{j}{1}$
  por definici\'on. Para $n\ge2$, supuesto $(\Pmat^{n-1})_{ij}=\pijn{i}{j}{n-1}$, por
  Chapman-Kolmogorov:
  \[
    (\Pmat^n)_{ij}=\sum_k(\Pmat^{n-1})_{ik}\pij{k}{j}
    =\sum_k\pijn{i}{k}{n-1}\pijn{k}{j}{1}=\pijn{i}{j}{n}.
  \]
  \textbf{(3)} Por (2): $(\boldsymbol{\mu}^{(0)}\Pmat^n)_j=\sum_i\mu_i^{(0)}(\Pmat^n)_{ij}
  =\sum_i\Prob{X_0=i}\,\pijn{i}{j}{n}=\Prob{X_n=j}=\mu_j^{(n)}$.\\[4pt]
  \textbf{(4)} $(\Pmat\mathbf{1})_i=\sum_j\pij{i}{j}=1$ para todo $i$, luego $\mathbf{1}$ es autovector derecho de $\lambda=1$.\\[4pt]
  \textbf{(5)} Sea $\Pmat\mathbf{v}=\lambda\mathbf{v}$ con $\|\mathbf{v}\|_\infty=1$ y $|v_{i^*}|=1$.
  Entonces $|\lambda|=|(\Pmat\mathbf{v})_{i^*}|=\bigl|\sum_j p_{i^*j}v_j\bigr|\le\sum_j p_{i^*j}|v_j|\le1$.
\end{demostracion}

\begin{nota}
  \textbf{Teorema de Perron-Frobenius (versi\'on estoc\'astica).}
  Para $\Pmat$ irreducible y aperi\'odica: $\lambda=1$ es el \'unico autovalor con $|\lambda|=1$
  y tiene multiplicidad algebraica $1$. Los dem\'as autovalores satisfacen $|\lambda_k|<1$,
  lo que controla la \emph{velocidad de convergencia} a la distribuci\'on estacionaria
  (ver \S4.6).
\end{nota}

% =============================================================================
\secheader{3}{Clasificaci\'on de Estados}

\subsecheader{3.1 \textperiodcentered\ Accesibilidad y comunicaci\'on}

\begin{definicion}{Accesibilidad}
  El estado $j$ es \textbf{accesible} desde $i$ (notaci\'on: $i \to j$) si existe $n \ge 0$
  tal que $\pijn{i}{j}{n} > 0$, es decir, es posible llegar de $i$ a $j$ en alg\'un n\'umero
  de pasos.\\[3pt]
  Dos estados $i$ y $j$ se \textbf{comunican} ($i \leftrightarrow j$) si $i \to j$ y $j \to i$.
  La comunicaci\'on es una relaci\'on de equivalencia en $\Sset$.
\end{definicion}

\begin{demostracion}
  Verificamos las tres propiedades de relaci\'on de equivalencia.\\[3pt]
  \textbf{Reflexividad} ($i\leftrightarrow i$): por convenio $\pijn{i}{i}{0}=1>0$, as\'i que $i\to i$.\\[3pt]
  \textbf{Simetr\'ia} ($i\leftrightarrow j\Rightarrow j\leftrightarrow i$): inmediata de la definici\'on
  (``$i\to j$ y $j\to i$'' es sim\'etrica en $i$ y $j$).\\[3pt]
  \textbf{Transitividad} ($i\leftrightarrow j$ y $j\leftrightarrow k\Rightarrow i\leftrightarrow k$):
  existen $r,s\ge0$ con $\pijn{i}{j}{r}>0$ y $\pijn{j}{k}{s}>0$.
  Por Chapman-Kolmogorov:
  $\pijn{i}{k}{r+s}\ge\pijn{i}{j}{r}\cdot\pijn{j}{k}{s}>0$, luego $i\to k$.
  An\'alogamente $k\to i$, de modo que $i\leftrightarrow k$. $\square$
\end{demostracion}

\begin{definicion}{Clase de comunicaci\'on}
  Las clases de equivalencia de $\leftrightarrow$ se llaman \textbf{clases de comunicaci\'on}.
  Una clase $C$ es \textbf{cerrada} si no existe salida: $i \in C$, $i \to j \Rightarrow j \in C$.\\[3pt]
  Una cadena es \textbf{irreducible} si solo tiene una clase de comunicaci\'on, es decir,
  todos los estados se comunican entre s\'i.
\end{definicion}

\begin{ejemplo}{Clases de comunicaci\'on en una cadena de 4 estados}
  Considerar la cadena con estados $\{1,2,3,4\}$ y la siguiente matriz de transici\'on:
  \[
    \Pmat = \begin{pmatrix}
      0   & 1   & 0   & 0   \\
      0.5 & 0.5 & 0   & 0   \\
      0.3 & 0   & 0.4 & 0.3 \\
      0   & 0   & 0   & 1
    \end{pmatrix}
  \]
  \begin{itemize}
    \item Estados $\{1,2\}$: $1 \leftrightarrow 2$ \quad (clase cerrada).
    \item Estado $\{3\}$: puede ir a $\{1\}$ y a $\{4\}$, pero ni $1$ ni $4$ van a $3$
          (clase abierta, estado transitorio).
    \item Estado $\{4\}$: $p_{44}=1$, no hay salida \quad (clase cerrada, estado absorbente).
  \end{itemize}
\end{ejemplo}

\begin{center}
\begin{tikzpicture}[>=Latex, node distance=2.5cm,
    state/.style={circle, draw=sectionbg, fill=thmbg, line width=1pt,
                  minimum size=0.9cm, font=\small\bfseries},
    abs/.style={circle, draw=warnborder, fill=warnbg, line width=1pt,
                minimum size=0.9cm, font=\small\bfseries},
    trans/.style={circle, draw=noteborder, fill=notebg, line width=1pt,
                  minimum size=0.9cm, font=\small\bfseries}]
  \node[state] (1) {$1$};
  \node[state, right of=1] (2) {$2$};
  \node[trans, below of=1, xshift=1.2cm] (3) {$3$};
  \node[abs, right of=3, xshift=0.5cm] (4) {$4$};
  \draw[->, thick, color=sectionbg]
    (1) edge[bend left=20] node[above,font=\tiny]{$1$} (2)
    (2) edge[bend left=20] node[below,font=\tiny]{$0.5$} (1)
    (2) edge[loop right] node[right,font=\tiny]{$0.5$} (2);
  \draw[->, thick, color=noteborder]
    (3) edge node[left,font=\tiny]{$0.3$} (1)
    (3) edge[loop left] node[left,font=\tiny]{$0.4$} (3)
    (3) edge node[below,font=\tiny]{$0.3$} (4);
  \draw[->, thick, color=warnborder]
    (4) edge[loop right] node[right,font=\tiny]{$1$} (4);
  \node[font=\tiny\itshape, color=sectionbg!70!black] at (2,-3.3)
    {Azul: clase cerrada $\{1,2\}$ \quad Naranja: transitorio $\{3\}$ \quad Rojo: absorbente $\{4\}$};
\end{tikzpicture}
\end{center}

\subsecheader{3.2 \textperiodcentered\ N\'umero de visitas a un estado}

\begin{definicion}{N\'umero de visitas}
  Dado un estado $i\in\Sset$, el \textbf{n\'umero total de visitas} a $i$ a lo largo de la
  trayectoria (contando desde $n=1$) es la variable aleatoria
  \[
    N(i) \;=\; \sum_{n=1}^{\infty} \ind_{\{X_n = i\}}.
  \]
  Si la cadena parte de $X_0=i$, entonces $N(i)$ cuenta cu\'antas veces el proceso
  \emph{regresa} a $i$.
\end{definicion}

\begin{proposicion}{Esperanza de visitas}
  \[
    \E{N(i)\mid X_0=i} \;=\; \sum_{n=1}^{\infty} \pijn{i}{i}{n}.
  \]
  En consecuencia: $i$ recurrente $\Longleftrightarrow$ $\E{N(i)\mid X_0=i}=\infty$.
\end{proposicion}

\begin{demostracion}
  Por linealidad y el teorema de Tonelli (sumandos no negativos):
  \[
    \E{N(i)\mid X_0=i}=\E{\sum_{n=1}^\infty\ind_{\{X_n=i\}}\;\Big|\;X_0=i}
    =\sum_{n=1}^\infty\Prob{X_n=i\mid X_0=i}=\sum_{n=1}^\infty\pijn{i}{i}{n}.\quad\square
  \]
\end{demostracion}

\begin{proposicion}{Distribuci\'on geom\'etrica de $N(i)$}
  Para todo $k\ge0$:
  \[
    \Prob{N(i)\ge k\mid X_0=i} = \fij{i}{i}^{\,k},
  \]
  donde $\fij{i}{i}=\Prob{\exists\,n\ge1: X_n=i\mid X_0=i}$ es la probabilidad de retorno.
  Es decir, $N(i)\mid\{X_0=i\}$ tiene distribuci\'on \textbf{Geom\'etrica} con par\'ametro
  $\fij{i}{i}$.
  \begin{itemize}
    \item Si $\fij{i}{i}<1$ (transitorio): $\E{N(i)\mid X_0=i}=\dfrac{\fij{i}{i}}{1-\fij{i}{i}}<\infty$.
    \item Si $\fij{i}{i}=1$ (recurrente): $\Prob{N(i)=\infty\mid X_0=i}=1$.
  \end{itemize}
\end{proposicion}

\begin{demostracion}
  Por la \textbf{propiedad de Markov fuerte} aplicada en el tiempo de parada $T_i^{(k)}$
  ($k$-\'esimo retorno a $i$): dado que la cadena regres\'o a $i$, el proceso
  reiniciado desde $i$ es independiente del pasado. As\'i:
  \[
    \Prob{N(i)\ge k\mid X_0=i}=\Prob{N(i)\ge k-1\mid X_0=i}\cdot\fij{i}{i},\quad k\ge1.
  \]
  Con $\Prob{N(i)\ge0}=1$, por inducci\'on: $\Prob{N(i)\ge k\mid X_0=i}=\fij{i}{i}^k$. $\square$
\end{demostracion}

\begin{nota}
  Combinando ambas proposiciones:
  \[
    i \text{ recurrente} \;\Longleftrightarrow\; \fij{i}{i}=1
    \;\Longleftrightarrow\; N(i)=\infty\;\text{c.s.}
    \;\Longleftrightarrow\; \sum_{n=1}^\infty\pijn{i}{i}{n}=\infty.
  \]
  Para un estado \emph{transitorio}, la cadena lo visita un n\'umero \emph{finito} de veces
  casi seguramente, y luego lo abandona para siempre.
\end{nota}

\subsecheader{3.3 \textperiodcentered\ Estados recurrentes y transitorios}

\begin{definicion}{Probabilidad de retorno}
  Define $f_{ii} = \Prob{X_n = i \text{ para alg\'un } n \ge 1 \mid X_0 = i}$, la
  probabilidad de que la cadena regrese al estado $i$. Entonces:
  \begin{itemize}
    \item $i$ es \textbf{recurrente} si $f_{ii} = 1$ (el retorno es seguro).
    \item $i$ es \textbf{transitorio} si $f_{ii} < 1$ (hay probabilidad positiva de no regresar).
  \end{itemize}
\end{definicion}

\begin{teorema}{Criterio de recurrencia por series}
  El estado $i$ es recurrente si y solo si
  \[
    \sum_{n=0}^{\infty} \pijn{i}{i}{n} = \infty.
  \]
  El estado $i$ es transitorio si y solo si $\displaystyle\sum_{n=0}^{\infty}\pijn{i}{i}{n} < \infty$,
  en cuyo caso el n\'umero esperado de visitas a $i$ es finito.
\end{teorema}

\begin{demostracion}
  Usamos los resultados de \S3.2.\\[3pt]
  \textbf{Paso 1.} Por la Proposici\'on de esperanza de visitas:
  $\E{N(i)\mid X_0=i}=\sum_{n=1}^\infty\pijn{i}{i}{n}$.\\[3pt]
  \textbf{Paso 2.} Por la Proposici\'on de distribuci\'on geom\'etrica: $N(i)\mid\{X_0=i\}\sim\mathrm{Geom}(\fij{i}{i})$, luego
  \[
    \E{N(i)\mid X_0=i}=\begin{cases}\dfrac{\fij{i}{i}}{1-\fij{i}{i}}<\infty & \text{si }\fij{i}{i}<1,\\[6pt]+\infty & \text{si }\fij{i}{i}=1.\end{cases}
  \]
  \textbf{Paso 3.} Uniendo: $\fij{i}{i}=1\Leftrightarrow\E{N(i)}=\infty\Leftrightarrow\sum_n\pijn{i}{i}{n}=\infty$.
  Esto es exactamente el enunciado. $\square$
\end{demostracion}

\begin{proposicion}{Recurrencia en cadenas finitas irreducibles}
  En una cadena de Markov con espacio de estados \emph{finito} e \emph{irreducible},
  todos los estados son \textbf{recurrentes positivos}.
\end{proposicion}

\begin{demostracion}
  Como $|\Sset| < \infty$, al menos un estado debe ser visitado infinitas veces (la cadena
  no puede ``escapar'' a infinito). En una cadena irreducible, si un estado $i$ es recurrente
  entonces \emph{todos} los estados son recurrentes (por comunicaci\'on). Y en $\Sset$ finito
  la recurrencia implica recurrencia positiva (el tiempo medio de retorno es finito).
\end{demostracion}

\begin{definicion}{Recurrencia positiva y nula}
  Sea $i$ recurrente y $\muu{i} = \E{\min\{n \ge 1 : X_n = i\} \mid X_0 = i}$ el
  \textbf{tiempo medio de retorno} a $i$. Entonces:
  \begin{itemize}
    \item $i$ es \textbf{recurrente positivo} si $\muu{i} < \infty$.
    \item $i$ es \textbf{recurrente nulo} si $\muu{i} = \infty$.
  \end{itemize}
\end{definicion}

\begin{cuidado}
  La recurrencia nula solo aparece en cadenas con espacio de estados \emph{infinito numerable}.
  El ejemplo m\'as conocido es la caminata aleatoria sim\'etrica en $\mathbb{Z}$: todos los
  estados son recurrentes (vuelve con probabilidad $1$), pero el tiempo esperado de retorno
  al origen es $\infty$. Ver Ejemplo~\ref{ex:paseo} (\S3.8) para el an\'alisis completo.
\end{cuidado}

\subsecheader{3.4 \textperiodcentered\ Recurrencia como propiedad de clase}

\begin{teorema}{Recurrencia se propaga por comunicaci\'on}
  Si $i\leftrightarrow j$, entonces $i$ recurrente $\Longleftrightarrow$ $j$ recurrente.
  Dentro de una clase de comunicaci\'on, \emph{todos} los estados son recurrentes o
  \emph{todos} son transitorios.
\end{teorema}

\begin{demostracion}
  Como $i\leftrightarrow j$, existen $r,s\ge1$ con $\alpha:=\pijn{i}{j}{r}>0$ y
  $\beta:=\pijn{j}{i}{s}>0$. Para todo $n\ge0$, por Chapman-Kolmogorov:
  \[
    \pijn{j}{j}{r+n+s} \;\ge\; \pijn{j}{i}{s}\cdot\pijn{i}{i}{n}\cdot\pijn{i}{j}{r}
    = \alpha\beta\cdot\pijn{i}{i}{n}.
  \]
  Sumando sobre $n$: si $i$ es recurrente ($\sum_n\pijn{i}{i}{n}=\infty$), entonces
  $\sum_n\pijn{j}{j}{n}=\infty$, luego $j$ recurrente. El rec\'iproco es sim\'etrico. $\square$
\end{demostracion}

\begin{proposicion}{Estados transitorios en clases abiertas}
  Si una clase de comunicaci\'on $C$ tiene salida (existe $i\in C$, $j\notin C$ con $\pij{i}{j}>0$),
  entonces todos los estados de $C$ son transitorios.
\end{proposicion}

\begin{demostracion}
  Partiendo de $i\in C$, con probabilidad $\pij{i}{j}>0$ la cadena alcanza $j\notin C$.
  Desde $j$, el retorno a $i$ puede ser imposible o tener probabilidad $<1$,
  luego $\fij{i}{i}<1$: $i$ es transitorio. Por el teorema anterior, toda la clase $C$ es transitoria. $\square$
\end{demostracion}

\subsecheader{3.5 \textperiodcentered\ Clasificaci\'on completa de estados}

\begin{teorema}{Clasificaci\'on tripartita}
  En una cadena de Markov irreducible, exactamente uno de los tres casos:
  \begin{enumerate}
    \item \textbf{Transitoria:} todos los estados son transitorios ($\fij{i}{i}<1$, $\E{N(i)}<\infty$).
    \item \textbf{Recurrente nula:} todos recurrentes con $\muu{i}=\infty$ para todo $i$.
    \item \textbf{Recurrente positiva:} todos recurrentes con $\muu{i}<\infty$ para todo $i$.
  \end{enumerate}
  Adem\'as: la cadena es recurrente positiva $\Longleftrightarrow$ admite distribuci\'on
  estacionaria $\boldsymbol{\pi}$ con $\pi_i>0$. En ese caso $\pi_i=1/\muu{i}$.
\end{teorema}

\begin{demostracion}
  La clasificaci\'on se propaga en la clase irreducible por el teorema anterior.
  Que todos tengan $\muu{i}<\infty$ o todos $\muu{i}=\infty$ se sigue de la cota
  $\muu{j}\le(1/\alpha\beta)\muu{i}$ obtenida de la desigualdad
  $\pijn{j}{j}{r+n+s}\ge\alpha\beta\,\pijn{i}{i}{n}$.\\[4pt]
  \textbf{Lema de Kac (positivo $\Rightarrow$ existe $\boldsymbol{\pi}$).}
  Para $i$ recurrente positivo, definir
  $\pi_j = \E{\#\{0\le n < T_i : X_n=j\}\mid X_0=i}/\muu{i}$.
  Se puede verificar que $\boldsymbol{\pi}\Pmat=\boldsymbol{\pi}$ y $\sum_j\pi_j=1$,
  con $\pi_j=1/\muu{j}>0$ (por irreducibilidad). $\square$\\[3pt]
  \textbf{Rec\'iproco.} Si $\exists\,\boldsymbol{\pi}$ estacionaria con $\pi_i>0$,
  el Teorema Erg\'odico (\S4.3) implica que $X_n=i$ para infinitos $n$ (recurrencia)
  y la frecuencia de visita es $\pi_i>0$, luego $\muu{i}=1/\pi_i<\infty$. $\square$
\end{demostracion}

\begin{proposicion}{Finito e irreducible implica recurrencia positiva}
  Toda cadena con $|\Sset|<\infty$ irreducible es recurrente positiva.
\end{proposicion}

\begin{demostracion}
  En espacio finito la cadena no puede escapar al infinito, as\'i que al menos un estado
  es recurrente. Por irreducibilidad, todos lo son. Si alguno fuera recurrente nulo,
  $\pi_i=1/\muu{i}=0$ para todo $i$, por lo que $\sum_i\pi_i=0\ne1$: contradicci\'on. $\square$
\end{demostracion}

\subsecheader{3.6 \textperiodcentered\ Periodicidad}

\begin{definicion}{Periodo de un estado}
  El \textbf{periodo} del estado $i$ es el m\'aximo com\'un divisor de los tiempos en los
  que es posible regresar a $i$:
  \[
    d(i) = \gcd\!\bigl\{n \ge 1 : \pijn{i}{i}{n} > 0\bigr\}.
  \]
  Si $d(i) = 1$, el estado es \textbf{aperi\'odico}. Si $d(i) > 1$, es \textbf{peri\'odico}
  con periodo $d(i)$.
\end{definicion}

\begin{proposicion}{Periodicidad es invariante en una clase}
  Si $i \leftrightarrow j$, entonces $d(i) = d(j)$.
  Por tanto, el periodo es una propiedad de la clase de comunicaci\'on, no del estado individual.
\end{proposicion}

\begin{demostracion}
  Sean $r,s\ge1$ con $\pijn{i}{j}{r}>0$ y $\pijn{j}{i}{s}>0$.\\[3pt]
  \textbf{Paso 1: $d(j)\mid d(i)$.}
  Para cualquier $n$ con $\pijn{i}{i}{n}>0$, por Chapman-Kolmogorov:
  \[
    \pijn{j}{j}{r+n+s}\ge\pijn{j}{i}{s}\cdot\pijn{i}{i}{n}\cdot\pijn{i}{j}{r}>0,
  \]
  luego $r+n+s\in\{m:\pijn{j}{j}{m}>0\}$ y por tanto $d(j)\mid(r+n+s)$.
  Tomando dos valores $n_1,n_2$ con $\pijn{i}{i}{n_k}>0$: $d(j)\mid(r+n_1+s)$ y $d(j)\mid(r+n_2+s)$,
  de donde $d(j)\mid(n_1-n_2)$.
  Como $d(i)=\gcd\{n:\pijn{i}{i}{n}>0\}$ es la m\'inima combinaci\'on lineal entera positiva
  de ese conjunto, y $d(j)$ divide todas las diferencias de sus elementos: $d(j)\mid d(i)$.\\[3pt]
  \textbf{Paso 2: $d(i)\mid d(j)$.}
  Mismo argumento intercambiando $i$ y $j$ (con $r'=s$, $s'=r$).\\[3pt]
  \textbf{Conclusi\'on:} $d(j)\mid d(i)$ y $d(i)\mid d(j)$ implican $d(i)=d(j)$. $\square$
\end{demostracion}

\begin{ejemplo}{Periodicidad}
  \begin{itemize}
    \item Cadena $1 \to 2 \to 1$ (dos estados en ciclo): $d(1) = d(2) = 2$ (peri\'odica).
    \item Cadena del clima (con autoloops): $d(S) = d(L) = 1$ (aperi\'odica).
    \item Caminata aleatoria en $\mathbb{Z}$ (pasos $\pm 1$): $d(0) = 2$ (peri\'odica). Ver \S3.8.
  \end{itemize}
\end{ejemplo}

\begin{proposicion}{Convergencia en m\'ultiplos del per\'iodo}
  Sea $\{X_n\}$ irreducible con per\'iodo $d$ y distribuci\'on estacionaria $\boldsymbol{\pi}$.
  Para todo $i,j\in\Sset$:
  \[
    \lim_{n\to\infty} \pijn{i}{j}{nd} = d\cdot\pi_j.
  \]
  Para cadenas peri\'odicas no existe el l\'imite $\pijn{i}{j}{n}$, pero la \textbf{media de Ces\`aro} converge:
  \[
    \frac{1}{N}\sum_{n=0}^{N-1}\pijn{i}{j}{n} \;\xrightarrow{N\to\infty}\; \pi_j.
  \]
\end{proposicion}

\begin{demostracion}
  La convergencia de Ces\`aro es consecuencia directa del Teorema Erg\'odico (\S4.3):
  $\frac{1}{N}\sum_{n=0}^{N-1}\pijn{i}{j}{n}=\E{\frac{1}{N}\sum_{n=0}^{N-1}\ind_{\{X_n=j\}}\mid X_0=i}\to\pi_j$
  por convergencia dominada (el integrando est\'a acotado por $1$). $\square$
\end{demostracion}

\subsecheader{3.7 \textperiodcentered\ Cadenas erg\'odicas}

\begin{definicion}{Cadena erg\'odica}
  Una cadena de Markov es \textbf{erg\'odica} si es:
  \begin{enumerate}
    \item \textbf{Irreducible}: todos los estados se comunican.
    \item \textbf{Recurrente positiva}: todos los estados tienen tiempo medio de retorno finito.
    \item \textbf{Aperi\'odica}: $d(i) = 1$ para todo $i \in \Sset$.
  \end{enumerate}
  En particular, toda cadena irreducible con $\Sset$ finito y aperi\'odica es erg\'odica.
\end{definicion}

\begin{nota}
  La erg\'odicidad es la condici\'on que garantiza la convergencia de $\Pmat^n$ a una
  distribuci\'on estacionaria \'unica (ver Secci\'on~4).
  En $\Sset$ finito, irreducibilidad + aperiodicidad $\Rightarrow$ erg\'odicidad.
\end{nota}

\subsecheader{3.8 \textperiodcentered\ Ejemplo: Paseo aleatorio sim\'etrico en $\mathbb{Z}$}
\label{ex:paseo}

\begin{ejemplo}{Paseo aleatorio sim\'etrico en $\mathbb{Z}$}
  \textbf{Definici\'on.}
  Sean $\xi_1,\xi_2,\ldots$ i.i.d.\ con $\Prob{\xi_k=+1}=\Prob{\xi_k=-1}=\tfrac{1}{2}$.
  Define $X_0=0$, $X_n=\xi_1+\cdots+\xi_n$. El espacio de estados es $\mathbb{Z}$ con
  $\pij{i}{i+1}=\pij{i}{i-1}=\tfrac{1}{2}$ (y $0$ en otro caso).

  \begin{center}
  \begin{tikzpicture}[>=Latex, node distance=1.8cm,
      state/.style={circle, draw=sectionbg, fill=thmbg, line width=0.8pt,
                    minimum size=0.7cm, font=\small\bfseries}]
    \foreach \x in {-2,-1,0,1,2}
      \node[state] (n\x) at (\x*1.8,0) {$\x$};
    \foreach \x [count=\xi] in {-2,-1,0,1} {
      \pgfmathtruncatemacro{\xn}{\x+1}
      \draw[->, thick, color=sectionbg, bend left=22]
        (n\x) edge node[above,font=\tiny]{$\frac{1}{2}$} (n\xn);
      \draw[->, thick, color=defborder, bend left=22]
        (n\xn) edge node[below,font=\tiny]{$\frac{1}{2}$} (n\x);
    }
    \node[left=0.2cm of n-2, font=\normalsize\itshape, color=gray] {$\cdots$};
    \node[right=0.2cm of n2, font=\normalsize\itshape, color=gray] {$\cdots$};
    \node[font=\small\itshape, color=sectionbg!70!black] at (0,-1.1)
      {Paseo aleatorio sim\'etrico en $\mathbb{Z}$};
  \end{tikzpicture}
  \end{center}

  \textbf{Paso 1: Irreducibilidad.}
  Para $j>i$: $j$ es accesible desde $i$ en $j-i$ pasos hacia la derecha.
  Por simetr\'ia $i$ es accesible desde $j$. Luego $i\leftrightarrow j$ para todo par.

  \textbf{Paso 2: Per\'iodo $d(0)=2$.}
  Solo se puede regresar al origen en n\'umero \emph{par} de pasos (se necesitan
  exactamente tantos $+1$ como $-1$). As\'i $\pijn{0}{0}{2k+1}=0$ para todo $k$, y
  $\pijn{0}{0}{2k}>0$ para $k\ge1$. Luego $d(0)=\gcd\{2,4,6,\ldots\}=2$.

  \textbf{Paso 3: F\'ormula exacta para $\pijn{0}{0}{2n}$.}
  Regresar en $2n$ pasos requiere exactamente $n$ pasos $+1$ y $n$ pasos $-1$:
  \[
    \pijn{0}{0}{2n} = \binom{2n}{n}\frac{1}{4^n}.
  \]

  \textbf{Paso 4: Aproximaci\'on de Stirling ($n!\approx\sqrt{2\pi n}(n/e)^n$).}
  \[
    \binom{2n}{n}\frac{1}{4^n}
    = \frac{(2n)!}{(n!)^2\,4^n}
    \approx \frac{\sqrt{4\pi n}\,(2n/e)^{2n}}{2\pi n\,(n/e)^{2n}\,4^n}
    = \frac{\sqrt{4\pi n}}{2\pi n} = \frac{1}{\sqrt{\pi n}}.
  \]
  M\'as precisamente: $\binom{2n}{n}/4^n\sim 1/\sqrt{\pi n}$ cuando $n\to\infty$.

  \textbf{Paso 5: Recurrencia del estado $0$.}
  \[
    \sum_{n=1}^\infty \pijn{0}{0}{2n}
    = \sum_{n=1}^\infty\frac{\binom{2n}{n}}{4^n}
    \;\sim\;\sum_{n=1}^\infty\frac{1}{\sqrt{\pi n}} = \infty,
  \]
  pues $\sum n^{-1/2}$ diverge (criterio integral: $\int_1^\infty x^{-1/2}\,dx=\infty$).
  Por el criterio de recurrencia: estado $0$ es recurrente.
  Por irreducibilidad: \emph{todos} los estados son recurrentes.

  \textbf{Paso 6: Recurrencia nula.}
  Usando la identidad $\sum_{n=0}^\infty\binom{2n}{n}x^n=1/\sqrt{1-4x}$ ($|x|<1/4$):
  \[
    P_{00}(z) := \sum_{n=0}^\infty\pijn{0}{0}{2n}z^n
    = \sum_{n=0}^\infty\binom{2n}{n}\Bigl(\frac{z}{4}\Bigr)^n
    = \frac{1}{\sqrt{1-z}}, \quad |z|<1.
  \]
  Por la identidad $P_{00}(z)=1/(1-F_{00}(z))$ (\S3.9):
  $F_{00}(z)=1-\sqrt{1-z}\xrightarrow{z\to1^-}1$, confirmando $f_{00}=1$ (recurrente).
  Pero
  \[
    \muu{0} = F'_{00}(1^-) = \lim_{z\to1^-}\frac{1}{2\sqrt{1-z}} = +\infty.
  \]
  El estado $0$ (y por irreducibilidad, todos) son \textbf{recurrentes nulos}.

  \textbf{Conclusi\'on.}
  \begin{itemize}
    \item \textbf{Irreducible}, \textbf{peri\'odico} ($d=2$), \textbf{recurrente nulo}.
    \item Vuelve a cualquier estado con probabilidad $1$, pero en tiempo esperado infinito.
    \item \textbf{Sin distribuci\'on estacionaria:} $\pi_i=1/\muu{i}=0$ y no es normalizable.
  \end{itemize}
\end{ejemplo}

\begin{cuidado}
  Recurrente no significa regreso \emph{r\'apido}. El paseo sim\'etrico en $\mathbb{Z}$
  vuelve al origen c.s., pero $\muu{0}=\infty$. En contraste, toda cadena irreducible
  en espacio finito es recurrente \emph{positiva} con tiempos de retorno finitos.
\end{cuidado}

\subsecheader{3.9 \textperiodcentered\ Funciones generatrices para tiempos de primer paso}

\begin{definicion}{Funciones generatrices de transici\'on y de primer paso}
  Para $i,j\in\Sset$ y $|z|\le1$:
  \begin{align*}
    P_{ij}(z) &= \sum_{n=0}^{\infty}\pijn{i}{j}{n}z^n
    \quad\text{(generatriz de transici\'on),}\\[4pt]
    F_{ij}(z) &= \sum_{n=1}^{\infty}\fij{i}{j}^{(n)}z^n
    \quad\text{(generatriz de primer paso).}
  \end{align*}
  Evaluando en $z=1$: $F_{ij}(1)=\fij{i}{j}$ (prob.\ total de alcanzar $j$ desde $i$).
\end{definicion}

\begin{proposicion}{Identidades entre $P_{ij}$ y $F_{ij}$}
  \begin{align}
    P_{ii}(z) &= \frac{1}{1-F_{ii}(z)}, \label{eq:Pii-Fii}\\[4pt]
    P_{ij}(z) &= F_{ij}(z)\cdot P_{jj}(z), \quad i\ne j. \label{eq:Pij-Fij}
  \end{align}
\end{proposicion}

\begin{demostracion}
  \textbf{Caso diagonal.} Condicionando sobre el primer retorno a $i$ (instante $k$):
  \[
    \pijn{i}{i}{n}=\sum_{k=1}^n\fij{i}{i}^{(k)}\pijn{i}{i}{n-k},\quad n\ge1;\qquad\pijn{i}{i}{0}=1.
  \]
  Multiplicando por $z^n$ y sumando: $P_{ii}(z)=1+F_{ii}(z)P_{ii}(z)$, de donde \eqref{eq:Pii-Fii}.
  An\'alogamente para \eqref{eq:Pij-Fij}. $\square$
\end{demostracion}

\begin{proposicion}{Criterio de recurrencia v\'ia generatrices}
  \begin{itemize}
    \item $i$ recurrente $\Longleftrightarrow$ $F_{ii}(1)=1$ $\Longleftrightarrow$ $P_{ii}(z)\to\infty$ cuando $z\to1^-$.
    \item $i$ recurrente positivo $\Longleftrightarrow$ $\muu{i}=F'_{ii}(1^-)<\infty$. Entonces $\pi_i=1/\muu{i}$.
    \item $i$ recurrente nulo $\Longleftrightarrow$ $F_{ii}(1)=1$ pero $F'_{ii}(1^-)=+\infty$.
  \end{itemize}
\end{proposicion}

\begin{demostracion}
  De la identidad $P_{ii}(z)=1/(1-F_{ii}(z))$:
  \begin{itemize}
    \item Si $F_{ii}(1)<1$: $(1-F_{ii}(1))>0$, luego $P_{ii}(1)=1/(1-F_{ii}(1))<\infty$, es decir,
          $\sum_n\pijn{i}{i}{n}<\infty$: estado transitorio.
    \item Si $F_{ii}(1)=1$: el denominador $\to0$ cuando $z\to1^-$, luego $P_{ii}(z)\to+\infty$:
          $\sum_n\pijn{i}{i}{n}=\infty$: estado recurrente.
    \item Para distinguir positivo/nulo: derivando $P_{ii}(z)=1/(1-F_{ii}(z))$,
          $P'_{ii}(z)=F'_{ii}(z)/(1-F_{ii}(z))^2$. El tiempo medio de retorno es
          $\muu{i}=\sum_n n\,\fij{i}{i}^{(n)}=F'_{ii}(1^-)$. Si $\muu{i}<\infty$: positivo;
          si $\muu{i}=\infty$: nulo. $\square$
  \end{itemize}
\end{demostracion}

\begin{nota}
  La relaci\'on $P_{ii}(z)=1/(1-F_{ii}(z))$ es el puente algebraico entre las
  \emph{probabilidades de transici\'on} (directas de calcular v\'ia $\Pmat^n$) y
  las \emph{probabilidades de primer paso} (que caracterizan recurrencia).
  Se usa de forma central en el an\'alisis del paseo aleatorio (\S3.8) y en la
  teor\'ia de renovaci\'on.
\end{nota}

% =============================================================================
\secheader{4}{Distribuciones Estacionarias y Comportamiento a Largo Plazo}

\subsecheader{4.1 \textperiodcentered\ Distribuci\'on estacionaria}

\begin{definicion}{Distribuci\'on estacionaria (invariante)}
  Un vector de probabilidades $\boldsymbol{\pi} = (\pi_j)_{j \in \Sset}$ es una
  \textbf{distribuci\'on estacionaria} (o \emph{invariante}) si satisface:
  \begin{align}
    \boldsymbol{\pi}\,\Pmat &= \boldsymbol{\pi} \label{eq:stat1}\\
    \sum_{j \in \Sset} \pi_j &= 1,\quad \pi_j \ge 0 \quad \forall j. \label{eq:stat2}
  \end{align}
  En forma escalar: $\pi_j = \sum_{i \in \Sset} \pi_i\, \pij{i}{j}$ para todo $j$.\\[3pt]
  Si $\boldsymbol{\mu}^{(0)} = \boldsymbol{\pi}$, entonces $\boldsymbol{\mu}^{(n)} = \boldsymbol{\pi}$
  para todo $n$: la distribuci\'on no cambia en el tiempo.
\end{definicion}

\begin{nota}
  La distribuci\'on estacionaria existe siempre en una cadena \emph{irreducible recurrente
  positiva}. La unicidad se obtiene adem\'as si la cadena es irreducible.
\end{nota}

\subsecheader{4.2 \textperiodcentered\ C\'alculo de la distribuci\'on estacionaria}

\begin{ejemplo}{Distribuci\'on estacionaria del clima}
  Con $\Pmat = \bigl(\begin{smallmatrix}0.8 & 0.2 \\ 0.4 & 0.6\end{smallmatrix}\bigr)$,
  resolver $\boldsymbol{\pi}\Pmat = \boldsymbol{\pi}$, $\pi_S + \pi_L = 1$:
  \begin{align*}
    \pi_S &= 0.8\,\pi_S + 0.4\,\pi_L \\
    \pi_L &= 0.2\,\pi_S + 0.6\,\pi_L
  \end{align*}
  De la primera ecuaci\'on: $0.2\,\pi_S = 0.4\,\pi_L$, luego $\pi_S = 2\pi_L$.\\
  Con $\pi_S + \pi_L = 1$: $\pi_L = 1/3$, $\pi_S = 2/3$.\\[3pt]
  \textbf{Interpretaci\'on:} a largo plazo, $2/3$ de los d\'ias son soleados.
\end{ejemplo}

\begin{ejemplo}{Distribuci\'on estacionaria de una cadena de 3 estados}
  Sea $\Sset = \{1,2,3\}$ con:
  \[
    \Pmat = \begin{pmatrix}0 & 0.5 & 0.5 \\ 0.3 & 0.4 & 0.3 \\ 0.6 & 0 & 0.4\end{pmatrix}.
  \]
  Las ecuaciones $\boldsymbol{\pi}\Pmat = \boldsymbol{\pi}$ con $\pi_1+\pi_2+\pi_3=1$ dan:
  \begin{align*}
    \pi_1 &= 0.3\pi_2 + 0.6\pi_3 \\
    \pi_2 &= 0.5\pi_1 + 0.4\pi_2 \;\Rightarrow\; 0.6\pi_2 = 0.5\pi_1 \\
    \pi_3 &= 0.5\pi_1 + 0.3\pi_2 + 0.4\pi_3 \;\Rightarrow\; 0.6\pi_3 = 0.5\pi_1 + 0.3\pi_2
  \end{align*}
  Resolviendo: $\pi_2 = \tfrac{5}{6}\pi_1$, $\pi_3 = \tfrac{5}{6}\pi_1 - \tfrac{1}{2}\pi_2 = \tfrac{5}{12}\pi_1$.\\
  Normalizando $\pi_1(1 + 5/6 + 5/12)=1$: $\pi_1 = 12/27$, $\pi_2 = 10/27$, $\pi_3 = 5/27$.
\end{ejemplo}

\subsecheader{4.3 \textperiodcentered\ Teorema erg\'odico y convergencia}

\begin{teorema}{Convergencia para cadenas erg\'odicas}
  Sea $\{X_n\}$ una cadena de Markov erg\'odica (irreducible, recurrente positiva, aperi\'odica)
  con distribuci\'on estacionaria \'unica $\boldsymbol{\pi}$. Entonces para todo $i, j \in \Sset$:
  \[
    \lim_{n \to \infty} \pijn{i}{j}{n} = \pi_j,
  \]
  independientemente del estado inicial $i$. En notaci\'on matricial:
  $\Pmat^n \xrightarrow{n\to\infty} \boldsymbol{\Pi}$, donde $\boldsymbol{\Pi}$ es la
  matriz cuyas filas son todas iguales a $\boldsymbol{\pi}$.
\end{teorema}

\begin{demostracion}
  \textbf{Descomposici\'on espectral (caso $|\Sset|<\infty$).}\\[3pt]
  \textbf{Paso 1.} Por la cota espectral (\S2.6, propiedad 5), los autovalores $\lambda_k$ de $\Pmat$
  satisfacen $|\lambda_k|\le1$. Para una cadena erg\'odica (irreducible + aperi\'odica + finita),
  el Teorema de Perron-Frobenius garantiza que $\lambda_1=1$ es simple y $|\lambda_k|<1$ para $k\ge2$.\\[3pt]
  \textbf{Paso 2.} Si $\Pmat$ es diagonalizable con $\Pmat=\sum_{k=1}^s\lambda_k\mathbf{v}_k\mathbf{u}_k^\top$
  (donde $\mathbf{v}_k$, $\mathbf{u}_k$ son vectores propios derechos e izquierdos), entonces:
  \[
    \Pmat^n = \sum_{k=1}^s \lambda_k^n\,\mathbf{v}_k\mathbf{u}_k^\top.
  \]
  \textbf{Paso 3.} El t\'ermino $k=1$ da $\lambda_1^n=1$ con $\mathbf{v}_1=\mathbf{1}$ y
  $\mathbf{u}_1^\top=\boldsymbol{\pi}$, produciendo $\mathbf{1}\boldsymbol{\pi}=\boldsymbol{\Pi}$.
  Los t\'erminos $k\ge2$ tienen $|\lambda_k^n|=|\lambda_k|^n\to0$, luego:
  \[
    \Pmat^n - \boldsymbol{\Pi} = \sum_{k=2}^s\lambda_k^n\,\mathbf{v}_k\mathbf{u}_k^\top
    \xrightarrow{n\to\infty} 0,
  \]
  es decir, $\pijn{i}{j}{n}=(\Pmat^n)_{ij}\to\pi_j$ para todo par $i,j$. $\square$
\end{demostracion}

\begin{nota}
  Este teorema garantiza que, sin importar d\'onde comience la cadena, la distribuci\'on
  marginal converge a $\boldsymbol{\pi}$.
  Si la cadena es peri\'odica, el l\'imite $\pijn{i}{j}{n}$ puede no existir, pero
  la media Ces\`aro converge: $\frac{1}{N}\sum_{n=0}^{N-1}\pijn{i}{j}{n} \to \pi_j$.
\end{nota}

\begin{teorema}{Teorema erg\'odico (frecuencias de visita)}
  Sea $\{X_n\}$ irreducible y recurrente positiva. Para todo $j \in \Sset$ y toda
  trayectoria (con probabilidad 1):
  \[
    \frac{1}{N}\sum_{n=0}^{N-1} \ind_{\{X_n = j\}} \;\xrightarrow{N\to\infty}\; \pi_j.
  \]
  La proporci\'on de tiempo que la cadena pasa en el estado $j$ converge (c.s.) a $\pi_j$.
  Adem\'as, $\pi_j = 1/\muu{j}$ (rec\'iproco del tiempo medio de retorno).
\end{teorema}

\begin{demostracion}
  Demostraci\'on completa en \S4.5 via la Ley Fuerte de los Grandes N\'umeros.
  Esquema: los tiempos entre retornos a $j$ son i.i.d.\ con media $\muu{j}$
  (propiedad de Markov fuerte), luego la LFGN da $V_j(N)/N\to1/\muu{j}$ c.s.,
  que coincide con $\pi_j$ por la unicidad de la distribuci\'on estacionaria.
\end{demostracion}

\subsecheader{4.4 \textperiodcentered\ Reversibilidad y balance detallado}

\begin{definicion}{Condici\'on de balance detallado}
  Una distribuci\'on $\boldsymbol{\pi}$ satisface el \textbf{balance detallado} respecto a
  $\Pmat$ si
  \[
    \pi_i\,\pij{i}{j} = \pi_j\,\pij{j}{i} \quad \forall\, i, j \in \Sset.
  \]
  Si existe tal $\boldsymbol{\pi}$, la cadena es \textbf{reversible} (en equilibrio,
  no se puede distinguir si el tiempo avanza hacia adelante o hacia atr\'as).
\end{definicion}

\begin{proposicion}{Balance detallado implica estacionariedad}
  Si $\boldsymbol{\pi}$ satisface el balance detallado respecto a $\Pmat$ y
  $\sum_j \pi_j = 1$, entonces $\boldsymbol{\pi}$ es distribuci\'on estacionaria.
\end{proposicion}

\begin{demostracion}
  $\displaystyle\sum_{i} \pi_i\,\pij{i}{j} = \sum_{i} \pi_j\,\pij{j}{i}
  = \pi_j \sum_{i} \pij{j}{i} = \pi_j \cdot 1 = \pi_j$.
  Luego $\boldsymbol{\pi}\Pmat = \boldsymbol{\pi}$.
\end{demostracion}

\begin{cuidado}
  El rec\'iproco es falso: existen distribuciones estacionarias que no satisfacen
  balance detallado (cadenas no reversibles). Balance detallado es una condici\'on
  \emph{suficiente} para estacionariedad, no necesaria.
\end{cuidado}

\subsecheader{4.5 \textperiodcentered\ Demostraci\'on del Teorema Erg\'odico}

\begin{demostracion}
  (del Teorema de frecuencias de visita, \S4.3.)\\[4pt]
  Sea $\{X_n\}$ irreducible y recurrente positiva. Fijamos $j\in\Sset$ y sea
  $0<T_j^{(1)}<T_j^{(2)}<\cdots$ la sucesi\'on de tiempos de retorno a $j$ (partiendo de $X_0=j$).
  Por la \textbf{propiedad de Markov fuerte}, los tiempos entre retornos
  $\Delta_k=T_j^{(k)}-T_j^{(k-1)}$ son \textbf{i.i.d.}\ con media
  $\muu{j}=\E{\Delta_1}<\infty$.\\[4pt]
  Sea $V_j(N)=\#\{n\in\{0,\ldots,N-1\}:X_n=j\}$.
  La relaci\'on $T_j^{(V_j(N))}\le N-1<T_j^{(V_j(N)+1)}$ divide por $V_j(N)$:
  \[
    \frac{T_j^{(V_j(N))}}{V_j(N)} \;\le\; \frac{N-1}{V_j(N)} \;<\;
    \frac{T_j^{(V_j(N)+1)}}{V_j(N)}.
  \]
  Por la \textbf{Ley Fuerte de los Grandes N\'umeros} aplicada a $\{\Delta_k\}$:
  $T_j^{(k)}/k\to\muu{j}$ c.s., luego los dos extremos tienden a $\muu{j}$.
  Por tanto $N/V_j(N)\to\muu{j}$ c.s., es decir:
  \[
    \frac{V_j(N)}{N} = \frac{1}{N}\sum_{n=0}^{N-1}\ind_{\{X_n=j\}}
    \;\xrightarrow{N\to\infty}\; \frac{1}{\muu{j}} = \pi_j \quad\text{c.s.} \qquad\square
  \]
\end{demostracion}

\begin{nota}
  La propiedad de Markov fuerte convierte los tiempos de retorno en una sucesi\'on i.i.d.,
  permitiendo aplicar la LFGN. El Teorema Erg\'odico vale para cadenas peri\'odicas
  tambi\'en (no requiere aperiodicidad).
\end{nota}

\subsecheader{4.6 \textperiodcentered\ Velocidad de convergencia y brecha espectral}

\begin{teorema}{Convergencia exponencial para cadenas finitas erg\'odicas}
  Sea $\{X_n\}$ erg\'odica con $|\Sset|=s$ y autovalores $1=\lambda_1>|\lambda_2|\ge\cdots\ge|\lambda_s|$.
  Entonces existe $C>0$ tal que:
  \[
    \max_{i,j\in\Sset}\bigl|\pijn{i}{j}{n}-\pi_j\bigr| \;\le\; C\,\rho^n,
    \qquad \rho := |\lambda_2| < 1.
  \]
\end{teorema}

\begin{demostracion}
  Si $\Pmat$ es diagonalizable: $\Pmat=V\Lambda V^{-1}$ con $\Lambda=\mathrm{diag}(\lambda_1,\ldots,\lambda_s)$,
  luego $\Pmat^n=V\Lambda^n V^{-1}$.
  La componente de $\lambda_1=1$ da $\boldsymbol{\Pi}$ (l\'imite); los restantes $s-1$ t\'erminos
  decaen a tasa $\rho^n$. $\square$
\end{demostracion}

\begin{definicion}{Brecha espectral y tiempo de mezcla}
  La \textbf{brecha espectral} es $\gamma=1-\rho\in(0,1]$.
  El \textbf{tiempo de mezcla} con precisi\'on $\varepsilon>0$ es
  $t_{\mathrm{mix}}(\varepsilon)=\min\bigl\{n:\max_{i,j}|\pijn{i}{j}{n}-\pi_j|\le\varepsilon\bigr\}$.
  De la cota: $t_{\mathrm{mix}}(\varepsilon)\le\lceil\log(C/\varepsilon)/\log(1/\rho)\rceil$.
\end{definicion}

\begin{ejemplo}{Brecha espectral de la cadena del clima}
  $\Pmat=\bigl(\begin{smallmatrix}0.8&0.2\\0.4&0.6\end{smallmatrix}\bigr)$:
  $\operatorname{tr}(\Pmat)=1.4$, $\det(\Pmat)=0.40$.
  Autovalores: $\lambda_1=1$, $\lambda_2=0.4$. Brecha espectral $\gamma=0.6$.\\
  Convergencia: $|\pijn{i}{j}{n}-\pi_j|\le C\cdot0.4^n\to0$ exponencialmente.
  En $n=10$ pasos: $0.4^{10}\approx10^{-4}$, pr\'acticamente en r\'egimen estacionario.
\end{ejemplo}

\subsecheader{4.7 \textperiodcentered\ Algoritmo de Metropolis-Hastings}

\begin{nota}
  \textbf{Motivaci\'on (MCMC).} A menudo se desea muestrear de una distribuci\'on objetivo
  $\boldsymbol{\pi}$ difícil de evaluar directamente (por ejemplo en inferencia bayesiana).
  El truco: construir una cadena de Markov \emph{reversible} con $\boldsymbol{\pi}$ como
  estacionaria.
\end{nota}

\begin{definicion}{Algoritmo de Metropolis-Hastings}
  Dada una cadena \emph{propuesta} con $q_{ij}$ (cualquier matriz estoc\'astica) y objetivo
  $\boldsymbol{\pi}$, la cadena resultante tiene:
  \[
    p_{ij} = q_{ij}\cdot\min\!\left(1,\;\frac{\pi_j\,q_{ji}}{\pi_i\,q_{ij}}\right), \quad i\ne j;
    \qquad p_{ii} = 1-\sum_{k\ne i}p_{ik} \quad\text{(rechazo)}.
  \]
\end{definicion}

\begin{proposicion}{Correctitud: balance detallado respecto a $\boldsymbol{\pi}$}
  $\pi_i\,p_{ij}=\pi_j\,p_{ji}$ para todo $i\ne j$, luego $\boldsymbol{\pi}$ es estacionaria.
\end{proposicion}

\begin{demostracion}
  $\pi_i\,p_{ij}=\pi_i\,q_{ij}\cdot\min(1,\pi_j q_{ji}/\pi_i q_{ij})
  =\min(\pi_i q_{ij},\pi_j q_{ji})=\pi_j\,p_{ji}$
  (por simetr\'ia de $\min$). $\square$
\end{demostracion}

\begin{nota}
  Metropolis-Hastings es la base de los m\'etodos \textbf{MCMC} (Markov Chain Monte Carlo):
  permite simular distribuciones de alta dimensi\'on sin conocer la constante de normalizaci\'on,
  siendo fundamental en estad\'istica bayesiana, f\'isica estad\'istica y machine learning.
  El \textbf{muestreo de Gibbs} es el caso especial en que la tasa de aceptaci\'on es $1$
  (siempre se acepta la propuesta).
\end{nota}

% =============================================================================
\secheader{5}{Tiempos de Primer Paso y Absorci\'on}

\subsecheader{5.0 \textperiodcentered\ Funciones arm\'onicas y ecuaci\'on de Poisson}

\begin{definicion}{Funci\'on arm\'onica}
  Sea $A\subseteq\Sset$ (frontera). Una funci\'on $h:\Sset\to\mathbb{R}$ es
  \textbf{arm\'onica} en $\Sset\setminus A$ si:
  \[
    h(i)=\sum_{j\in\Sset}\pij{i}{j}\,h(j)=\E{h(X_1)\mid X_0=i}
    \quad\forall\,i\in\Sset\setminus A.
  \]
  Es decir, $h(i)$ es el valor esperado de $h$ en el siguiente paso.
  An\'alogamente, $h$ es \textbf{superarm\'onica} si $h(i)\ge\E{h(X_1)\mid X_0=i}$
  (el proceso $\{h(X_n)\}$ es supermartingala).
\end{definicion}

\begin{teorema}{Principio del m\'aximo discreto}
  Si $h$ es arm\'onica en todo $\Sset$ (sin frontera) y $|\Sset|<\infty$,
  entonces $h$ es constante en cada clase de comunicaci\'on cerrada.
\end{teorema}

\begin{demostracion}
  Sea $i^*=\arg\max_{i\in C}h(i)$ en una clase cerrada $C$.
  Como $h(i^*)=\sum_j\pij{i^*}{j}h(j)$ y $h(j)\le h(i^*)$ para $j\in C$
  (la clase es cerrada, luego $\pij{i^*}{j}=0$ si $j\notin C$), la igualdad
  requiere $h(j)=h(i^*)$ para todo $j$ con $\pij{i^*}{j}>0$.
  Por irreducibilidad de $C$, esto se extiende a toda la clase. $\square$
\end{demostracion}

\begin{proposicion}{Problemas de absorci\'on como problemas arm\'onicos}
  \begin{enumerate}
    \item Las \textbf{probabilidades de absorci\'on} $h_i^{(a)}=\Prob{X_{T_A}=a\mid X_0=i}$
          son arm\'onicas en $\Sset\setminus A$, con condiciones de frontera $h_a^{(a)}=1$
          y $h_b^{(a)}=0$ ($b\in A$, $b\ne a$).
    \item El \textbf{tiempo esperado de absorci\'on} $t_i=\E{T_A\mid X_0=i}$ satisface la
          \textbf{ecuaci\'on de Poisson} (arm\'onica inhomog\'enea):
          \[
            t_i = 1 + \sum_{j\notin A}\pij{i}{j}\,t_j, \quad i\notin A,
            \qquad t_a=0 \;\forall\,a\in A.
          \]
  \end{enumerate}
\end{proposicion}

\begin{demostracion}
  \textbf{(1)} Por la propiedad de Markov en el primer paso:
  $h_i^{(a)}=\sum_j\pij{i}{j}h_j^{(a)}$ (armonicidad).\\[3pt]
  \textbf{(2)} Condicionando en $X_1$:
  $t_i=1+\sum_j\pij{i}{j}\E{T_A\mid X_0=j}=1+\sum_{j\notin A}\pij{i}{j}t_j$
  (los t\'erminos $j\in A$ contribuyen $0$ porque $t_a=0$). $\square$
\end{demostracion}

\subsecheader{5.1 \textperiodcentered\ Tiempo de primer paso}

\begin{definicion}{Tiempo de primer paso (hitting time)}
  El \textbf{tiempo de primer paso} al estado $j$, partiendo de $i$, es:
  \[
    T_j = \min\{n \ge 1 : X_n = j\}, \quad T_j = \infty \text{ si nunca se llega a } j.
  \]
  Es un tiempo de parada respecto a la filtraci\'on natural.
  La \textbf{probabilidad de primer paso en $n$ pasos} de $i$ a $j$ es:
  \[
    \fij{i}{j}^{(n)} = \Prob{T_j = n \mid X_0 = i}, \quad
    \fij{i}{j} = \Prob{T_j < \infty \mid X_0 = i} = \sum_{n=1}^{\infty} \fij{i}{j}^{(n)}.
  \]
\end{definicion}

\begin{proposicion}{Relaci\'on entre $\pijn{i}{j}{n}$ y $\fij{i}{j}^{(n)}$}
  Para $n \ge 1$:
  \[
    \pijn{i}{j}{n} = \sum_{k=1}^{n} \fij{i}{j}^{(k)}\,\pijn{j}{j}{n-k}.
  \]
  Esto es una convolución: llegar de $i$ a $j$ en $n$ pasos = llegar por primera vez
  en $k$ pasos, y luego regresar de $j$ a $j$ en los $n-k$ restantes.
\end{proposicion}

\begin{demostracion}
  Particionamos el evento $\{X_n=j\}\cap\{X_0=i\}$ seg\'un el instante del primer arribo a $j$:
  \[
    \{X_n=j,\,X_0=i\} = \bigsqcup_{k=1}^{n}\{T_j=k,\,X_n=j,\,X_0=i\},
  \]
  donde $T_j=\min\{m\ge1:X_m=j\}$ (uniones disjuntas por el valor de $T_j$).
  Para cada $k\in\{1,\ldots,n\}$, condicionando en $\{T_j=k,X_0=i\}$ y aplicando la
  \textbf{propiedad de Markov fuerte} en el tiempo $T_j=k$:
  \[
    \Prob{X_n=j\mid T_j=k,\,X_0=i}=\Prob{X_{n-k}=j\mid X_0=j}=\pijn{j}{j}{n-k}.
  \]
  Luego, por probabilidad total:
  \[
    \pijn{i}{j}{n}=\Prob{X_n=j\mid X_0=i}
    =\sum_{k=1}^{n}\underbrace{\Prob{T_j=k\mid X_0=i}}_{=\fij{i}{j}^{(k)}}\cdot\pijn{j}{j}{n-k}.\quad\square
  \]
\end{demostracion}

\subsecheader{5.2 \textperiodcentered\ Tiempo medio de retorno y estacionariedad}

\begin{teorema}{Tiempo medio de retorno y distribuci\'on estacionaria}
  Sea $\{X_n\}$ irreducible y recurrente positiva, con distribuci\'on estacionaria
  $\boldsymbol{\pi}$. Entonces para todo $j \in \Sset$:
  \[
    \pi_j = \frac{1}{\muu{j}}, \quad \text{donde } \muu{j} = \E{T_j \mid X_0 = j}.
  \]
  Los estados m\'as frecuentemente visitados tienen tiempo medio de retorno m\'as corto.
\end{teorema}

\begin{demostracion}
  \textbf{Paso 1.} Fijamos $j$ y suponemos $X_0=j$.
  Por la \textbf{propiedad de Markov fuerte}, los tiempos entre retornos consecutivos
  $\Delta_k=T_j^{(k)}-T_j^{(k-1)}$ son i.i.d.\ con media $\muu{j}<\infty$.\\[3pt]
  \textbf{Paso 2.} Sea $V_j(N)=\#\{0\le n<N:X_n=j\}$. Como se demuestra en \S4.5,
  la Ley Fuerte de los Grandes N\'umeros aplicada a los $\Delta_k$ da:
  \[
    \frac{V_j(N)}{N}\xrightarrow{N\to\infty}\frac{1}{\muu{j}}\quad\text{c.s.}
  \]
  \textbf{Paso 3.} Por el Teorema Erg\'odico (\S4.3), esa misma frecuencia converge a $\pi_j$.
  Por unicidad del l\'imite: $\pi_j=1/\muu{j}$.\\[3pt]
  \textbf{Validez para $X_0\ne j$:} desde cualquier estado inicial, la cadena alcanza $j$
  en tiempo finito (recurrencia positiva) y la frecuencia a largo plazo es la misma. $\square$
\end{demostracion}

\begin{ejemplo}{Tiempo de retorno en la cadena del clima}
  Con $\pi_S = 2/3$ y $\pi_L = 1/3$:
  \[
    \muu{S} = \frac{1}{\pi_S} = \frac{3}{2} = 1.5 \text{ d\'ias}, \quad
    \muu{L} = \frac{1}{\pi_L} = 3 \text{ d\'ias}.
  \]
  En promedio, el tiempo entre dos d\'ias lluviosos consecutivos es de 3 d\'ias.
\end{ejemplo}

\subsecheader{5.3 \textperiodcentered\ Probabilidades de absorci\'on}

\begin{definicion}{Estado absorbente y probabilidad de absorci\'on}
  Un estado $a$ es \textbf{absorbente} si $\pij{a}{a} = 1$ (una vez alcanzado, no se escapa).
  Dado un conjunto de estados absorbentes $A \subset \Sset$ y un estado transitorio $i \notin A$,
  la \textbf{probabilidad de absorci\'on} en $a \in A$ partiendo de $i$ es:
  \[
    h_i^{(a)} = \Prob{X_{T_A} = a \mid X_0 = i},
  \]
  donde $T_A = \min\{n \ge 0 : X_n \in A\}$.
\end{definicion}

\begin{teorema}{Sistema de ecuaciones para probabilidades de absorci\'on}
  Las probabilidades $h_i^{(a)}$ satisfacen el sistema lineal:
  \[
    h_i^{(a)} = \pij{i}{a} + \sum_{k \notin A} \pij{i}{k}\,h_k^{(a)}, \quad i \notin A,
  \]
  con condici\'on de frontera $h_a^{(a)} = 1$ y $h_b^{(a)} = 0$ para $b \in A$, $b \neq a$.
\end{teorema}

\begin{demostracion}
  Condicionamos en el primer paso desde $i\notin A$:
  \begin{align*}
    h_i^{(a)}
    &= \Prob{X_{T_A}=a\mid X_0=i}\\[3pt]
    &= \sum_{j\in\Sset}\pij{i}{j}\,\Prob{X_{T_A}=a\mid X_1=j}
       \quad\text{(probabilidad total sobre el primer paso)}\\[3pt]
    &= \underbrace{\pij{i}{a}\cdot1}_{\text{si } j=a\in A}
       + \underbrace{\sum_{b\in A,\,b\ne a}\pij{i}{b}\cdot0}_{\text{si }j\in A,j\ne a}
       + \underbrace{\sum_{k\notin A}\pij{i}{k}\,h_k^{(a)}}_{\text{si }j\notin A}\\[3pt]
    &= \pij{i}{a} + \sum_{k\notin A}\pij{i}{k}\,h_k^{(a)}.
  \end{align*}
  El segundo t\'ermino usa la \textbf{propiedad de Markov}: si $X_1=k\notin A$,
  la probabilidad de absorci\'on en $a$ es la misma que partiendo de $k$. $\square$
\end{demostracion}

\subsecheader{5.4 \textperiodcentered\ Tiempo esperado de absorci\'on: matriz fundamental}

\begin{definicion}{Forma can\'onica y matriz fundamental}
  Ordenando los estados como (transitorios, absorbentes), la matriz de transici\'on toma
  la forma can\'onica:
  \[
    \Pmat = \begin{pmatrix} \mathbf{Q}_T & \mathbf{R} \\ \mathbf{0} & \mathbf{I} \end{pmatrix},
  \]
  donde $\mathbf{Q}_T$ es la submatriz de transici\'on entre estados transitorios y
  $\mathbf{R}$ las probabilidades de ir de transitorio a absorbente.\\[3pt]
  La \textbf{matriz fundamental} es $\mathbf{N} = (\mathbf{I} - \mathbf{Q}_T)^{-1}$.
  La entrada $n_{ij}$ = n\'umero esperado de visitas al estado transitorio $j$
  partiendo de $i$ (antes de la absorci\'on).
\end{definicion}

\begin{proposicion}{Tiempo esperado de absorci\'on}
  El vector de tiempos esperados de absorci\'on $\mathbf{t}$ (con $t_i = \E{T_A \mid X_0=i}$)
  satisface:
  \[
    \mathbf{t} = \mathbf{N}\,\mathbf{1} = (\mathbf{I} - \mathbf{Q}_T)^{-1}\mathbf{1},
  \]
  donde $\mathbf{1}$ es el vector de unos. Las probabilidades de absorci\'on en cada
  estado absorbente forman la matriz $\mathbf{B} = \mathbf{N}\,\mathbf{R}$.
\end{proposicion}

\begin{demostracion}
  \textbf{Paso 1: $\mathbf{N}=(\mathbf{I}-\mathbf{Q}_T)^{-1}$.}
  Sea $\mathbf{m}_{ij}=\E{V_j\mid X_0=i}$ el n\'umero esperado de visitas al estado
  transitorio $j$ antes de la absorci\'on. Condicionando en el primer paso:
  \[
    m_{ij}=\delta_{ij}+\sum_{k\notin A}(Q_T)_{ik}\,m_{kj}.
  \]
  En notaci\'on matricial: $\mathbf{M}=\mathbf{I}+\mathbf{Q}_T\mathbf{M}$, de donde
  $(\mathbf{I}-\mathbf{Q}_T)\mathbf{M}=\mathbf{I}$ y $\mathbf{M}=(\mathbf{I}-\mathbf{Q}_T)^{-1}=\mathbf{N}$.
  (La inversa existe porque $\mathbf{Q}_T^n\to\mathbf{0}$: todos los estados transitorios
  tienen absorci\'on casi segura.)\\[3pt]
  \textbf{Paso 2: $\mathbf{t}=\mathbf{N}\mathbf{1}$.}
  El tiempo de absorci\'on es la suma de las visitas a todos los estados transitorios:
  $T_A=\sum_{j\notin A}V_j$. Por linealidad:
  \[
    t_i=\E{T_A\mid X_0=i}=\sum_{j\notin A}\E{V_j\mid X_0=i}=\sum_j\mathbf{N}_{ij}=(\mathbf{N}\mathbf{1})_i.\quad\square
  \]
\end{demostracion}

\begin{ejemplo}{Juego de la ruina ($N=4$)}
  Un jugador tiene $k$ fichas ($k \in \{0,1,2,3,4\}$) y en cada ronda gana o pierde una
  con probabilidad $p=0.5$ y $1-p=0.5$. Los estados $\{0,4\}$ son absorbentes (ruina o victoria).
  Los estados transitorios son $\{1,2,3\}$.
  \[
    \mathbf{Q}_T = \begin{pmatrix}0 & 0.5 & 0 \\ 0.5 & 0 & 0.5 \\ 0 & 0.5 & 0\end{pmatrix},
    \quad
    \mathbf{I}-\mathbf{Q}_T = \begin{pmatrix}1 & -0.5 & 0 \\ -0.5 & 1 & -0.5 \\ 0 & -0.5 & 1\end{pmatrix}.
  \]
  Calculando la inversa:
  $\mathbf{N} = \frac{1}{4}\begin{pmatrix}3 & 2 & 1 \\ 2 & 4 & 2 \\ 1 & 2 & 3\end{pmatrix}$.\\[4pt]
  Tiempo esperado de absorci\'on: $\mathbf{t} = \mathbf{N}\mathbf{1}$:
  $t_1 = (3+2+1)/4 = 3/2$... \\
  Corrigiendo: $t_1 = 3$, $t_2 = 4$, $t_3 = 3$ (se verifica que desde el estado central
  se espera m\'as tiempo antes de arruinarse o ganar).\\[3pt]
  Probabilidad de ruina desde estado $k$: $h_k^{(0)} = (4-k)/4$.
  Desde estado~$2$: $h_2^{(0)} = 0.5$.
\end{ejemplo}

\begin{center}
\begin{tikzpicture}[>=Latex, node distance=2.0cm,
    abs/.style={circle, draw=warnborder, fill=warnbg, line width=1pt,
                minimum size=0.8cm, font=\small\bfseries},
    trans/.style={circle, draw=sectionbg, fill=thmbg, line width=1pt,
                  minimum size=0.8cm, font=\small\bfseries}]
  \node[abs] (0) {$0$};
  \node[trans, right of=0] (1) {$1$};
  \node[trans, right of=1] (2) {$2$};
  \node[trans, right of=2] (3) {$3$};
  \node[abs, right of=3] (4) {$4$};
  \draw[->, thick, color=sectionbg]
    (1) edge[bend left=25] node[above,font=\tiny]{$0.5$} (2)
    (2) edge[bend left=25] node[above,font=\tiny]{$0.5$} (3)
    (2) edge[bend left=25] node[below,font=\tiny]{$0.5$} (1)
    (3) edge[bend left=25] node[below,font=\tiny]{$0.5$} (2);
  \draw[->, thick, color=warnborder]
    (1) edge node[below,font=\tiny]{$0.5$} (0)
    (3) edge node[below,font=\tiny]{$0.5$} (4);
  \draw[->, thick, color=warnborder]
    (0) edge[loop left] node[left,font=\tiny]{$1$} (0)
    (4) edge[loop right] node[right,font=\tiny]{$1$} (4);
  \node[font=\small\itshape, color=sectionbg!70!black] at (4,-1.4)
    {Juego de la ruina con $N=4$: estados absorbentes $\{0,4\}$};
\end{tikzpicture}
\end{center}

\subsecheader{5.5 \textperiodcentered\ Juego de la ruina: caso general $p\ne\tfrac{1}{2}$}

\begin{ejemplo}{Ruina del apostador con probabilidad $p$ arbitraria}
  \textbf{Planteamiento.} Jugador con $k$ fichas ($1\le k\le N-1$), gana con prob.\ $p$,
  pierde con $q=1-p$. Sea $h_k=\Prob{\text{ruina}\mid X_0=k}$.
  Por armonicidad (\S5.0):
  \[
    h_k = p\,h_{k+1}+q\,h_{k-1}, \quad k=1,\ldots,N-1,
    \qquad h_0=1,\quad h_N=0.
  \]

  \textbf{Ecuaci\'on de diferencias.} La ecuaci\'on caracter\'istica $q r^2 - r + p=0$
  tiene ra\'ices $r=1$ y $r=\rho:=q/p$.

  \textbf{Caso $p=q=\tfrac{1}{2}$} (ra\'iz doble $r=1$): $h_k=A+Bk$.
  Con $h_0=1$, $h_N=0$:
  \[
    \boxed{h_k = \frac{N-k}{N}.}
  \]

  \textbf{Caso $p\ne q$} (ra\'ices distintas): $h_k=A+B\rho^k$.
  Con $A+B=1$ y $A+B\rho^N=0$:
  \[
    \boxed{h_k = \frac{\rho^k - \rho^N}{1-\rho^N},\qquad \rho=\frac{q}{p}.}
  \]

  \textbf{Tiempo esperado de absorci\'on.} La ecuaci\'on de Poisson da $t_k=1+p\,t_{k+1}+q\,t_{k-1}$,
  $t_0=t_N=0$. Soluciones:
  \[
    t_k = \begin{cases}
      k(N-k) & p=\tfrac{1}{2},\\[4pt]
      \dfrac{k}{q-p} - \dfrac{N}{q-p}\cdot\dfrac{1-\rho^k}{1-\rho^N} & p\ne\tfrac{1}{2}.
    \end{cases}
  \]

  \begin{center}
    \renewcommand{\arraystretch}{1.3}
    \begin{tabular}{@{}c c c@{}}
      \toprule
      $p$ & $h_5^{(N=10)}$ & $t_5^{(N=10)}$\\
      \midrule
      $0.4$ ($\rho=1.5$) & $\dfrac{1.5^5-1.5^{10}}{1-1.5^{10}}\approx 0.870$ & $\approx 19.5$ \\[4pt]
      $0.5$ & $0.500$ & $25.0$ \\[4pt]
      $0.6$ ($\rho=2/3$) & $\dfrac{(2/3)^5-(2/3)^{10}}{1-(2/3)^{10}}\approx 0.130$ & $\approx 19.5$ \\
      \bottomrule
    \end{tabular}
  \end{center}

  \textbf{L\'imite $N\to\infty$ (rival con capital infinito):}
  \[
    h_k \;\xrightarrow{N\to\infty}\;
    \begin{cases}
      1 & \text{si } p\le\tfrac{1}{2},\\
      \bigl(\tfrac{q}{p}\bigr)^k & \text{si } p>\tfrac{1}{2}.
    \end{cases}
  \]
  Si $p>1/2$: probabilidad positiva de sobrevivir indefinidamente. Si $p\le1/2$: la ruina es cierta.
\end{ejemplo}

\subsecheader{5.6 \textperiodcentered\ Modelo de Ehrenfest (difusi\'on entre urnas)}

\begin{ejemplo}{Modelo de Ehrenfest}
  \textbf{Planteamiento.} $N$ bolas en dos urnas $A$ y $B$. En cada paso se elige una
  bola al azar y se transfiere a la otra urna. Estado: $X_n=$ n\'umero de bolas en $A$.

  \textbf{Transiciones:} $\pij{k}{k+1}=(N-k)/N$ y $\pij{k}{k-1}=k/N$.
  Proceso de nacimiento y muerte discreto.

  \textbf{Irreducibilidad y per\'iodo.} Irreducible (todo $k$ alcanzable desde cualquier
  otro). Peri\'odico con $d=2$ (la paridad de $X_n$ cambia en cada paso, pues $\pij{k}{k}=0$).

  \textbf{Distribuci\'on estacionaria (balance detallado).}
  $\pi_k\,(N-k)/N=\pi_{k+1}\,(k+1)/N$ implica $\pi_{k+1}/\pi_k=(N-k)/(k+1)=\binom{N}{k+1}/\binom{N}{k}$.
  Iterando desde $\pi_0$:
  \[
    \pi_k=\binom{N}{k}\pi_0.
    \quad\text{Normalizando: }\sum_{k=0}^N\binom{N}{k}\pi_0=2^N\pi_0=1.
  \]
  \[
    \boxed{\pi_k=\binom{N}{k}\frac{1}{2^N}\;\sim\;\mathrm{Binomial}\!\left(N,\tfrac{1}{2}\right).}
  \]

  \textbf{Tiempo de retorno al estado $0$:}
  $\muu{0}=1/\pi_0=2^N$.
  Para $N=100$: $\muu{0}=2^{100}\approx10^{30}$ pasos.
  A $10^{10}$ pasos/seg: $\approx3\times10^{12}$ a\~nos de espera.

  \textbf{Brecha espectral.} Los autovalores de $\Pmat$ son $\lambda_k=1-2k/N$ ($k=0,\ldots,N$).
  Brecha espectral $\gamma=2/N$; tiempo de mezcla $t_{\mathrm{mix}}=O(N\log N)$.

  \textbf{Segunda ley de la termodin\'amica.} La cadena es recurrente positiva
  (espacio finito e irreducible), luego matem\'aticamente regresa al estado $0$.
  Pero $\muu{0}=2^N$ es \emph{inobservable} para $N\sim10^{23}$ (n\'umero de Avogadro).
  La irreversibilidad termodin\'amica no es imposibilidad matem\'atica, sino
  probabilidad astronómicamente peque\~na.
\end{ejemplo}

% =============================================================================
\secheader{6}{Cadenas de Markov en Tiempo Continuo (CMTC)}

\subsecheader{6.1 \textperiodcentered\ Definici\'on y propiedad de Markov continua}

\begin{definicion}{Cadena de Markov en tiempo continuo (CMTC)}
  Un proceso estoc\'astico $\{X(t)\}_{t\ge0}$ con espacio de estados discreto $\Sset$
  es una \textbf{cadena de Markov en tiempo continuo} si satisface la propiedad de Markov:
  \[
    \Prob{X(t+s) = j \mid X(s) = i,\, \{X(u)\}_{0\le u < s}}
    = \Prob{X(t+s)=j \mid X(s)=i}
  \]
  para todo $t, s \ge 0$ e $i,j \in \Sset$.
  La cadena es \textbf{homog\'enea} si la probabilidad de transici\'on
  $p_{ij}(t) = \Prob{X(t+s)=j \mid X(s)=i}$ no depende de $s$.
\end{definicion}

\subsecheader{6.2 \textperiodcentered\ Tiempos de permanencia}

\begin{teorema}{Distribuci\'on exponencial de los tiempos de permanencia}
  En una CMTC homog\'enea, el tiempo que el proceso permanece en el estado $i$ antes de saltar
  es una variable aleatoria $\tau_i \sim \mathrm{Exp}(q_i)$, donde $q_i \ge 0$ es la
  \textbf{tasa de salida} del estado $i$.\\[3pt]
  Esto se sigue de la propiedad de falta de memoria de la Exponencial: la \'unica
  distribuci\'on continua sin memoria es la Exponencial.
\end{teorema}

\begin{nota}
  Si $q_i = 0$, el estado $i$ es absorbente (la cadena nunca sale de él).
  El tiempo esperado en el estado $i$ es $\E{\tau_i} = 1/q_i$.
\end{nota}

\subsecheader{6.3 \textperiodcentered\ Generador infinitesimal $\generator$}

\begin{definicion}{Generador infinitesimal (Q-matrix)}
  La \textbf{Q-matrix} o \textbf{generador infinitesimal} $\generator = (q_{ij})$ tiene entradas:
  \begin{itemize}
    \item $q_{ij} \ge 0$ para $i \neq j$: \textbf{tasa de transici\'on} de $i$ a $j$.
    \item $q_{ii} = -\sum_{j \neq i} q_{ij} = -q_i \le 0$: la diagonal es la tasa de salida negativa.
  \end{itemize}
  Propiedades: cada fila de $\generator$ suma $0$, i.e., $\generator\,\mathbf{1} = \mathbf{0}$.\\[3pt]
  Relaci\'on con $\Pmat(t)$:
  \[
    \lim_{h \to 0^+} \frac{p_{ij}(h) - \delta_{ij}}{h} = q_{ij}.
  \]
\end{definicion}

\begin{ejemplo}{Q-matrix para una cadena de 2 estados}
  Sea $\Sset = \{0, 1\}$, tasa de ir de $0$ a $1$: $\lambda$, tasa de $1$ a $0$: $\mu$.
  \[
    \generator = \begin{pmatrix} -\lambda & \lambda \\ \mu & -\mu \end{pmatrix}.
  \]
\end{ejemplo}

\begin{center}
\begin{tikzpicture}[>=Latex, node distance=4cm,
    state/.style={circle, draw=sectionbg, fill=thmbg, line width=1pt,
                  minimum size=1.1cm, font=\small\bfseries}]
  \node[state] (0) {$0$};
  \node[state, right of=0] (1) {$1$};
  \draw[->, thick, color=sectionbg]
    (0) edge[bend left=25] node[above,font=\small]{$\lambda$} (1)
    (1) edge[bend left=25] node[below,font=\small]{$\mu$} (0);
  \node[font=\small\itshape, color=sectionbg!70!black] at (2,-1.2)
    {CMTC de 2 estados: tasas $\lambda$ y $\mu$};
\end{tikzpicture}
\end{center}

\subsecheader{6.4 \textperiodcentered\ Ecuaciones de Kolmogorov}

\begin{teorema}{Ecuaciones de Kolmogorov (progresiva y regresiva)}
  Sea $\mathbf{P}(t) = (p_{ij}(t))$ la matriz de probabilidades de transici\'on.
  Bajo condiciones de regularidad:
  \begin{align}
    \mathbf{P}'(t) &= \mathbf{P}(t)\,\generator \quad \text{(ecuaci\'on progresiva de Kolmogorov)}\\
    \mathbf{P}'(t) &= \generator\,\mathbf{P}(t) \quad \text{(ecuaci\'on regresiva de Kolmogorov)}
  \end{align}
  con condici\'on inicial $\mathbf{P}(0) = \mathbf{I}$.
  La soluci\'on formal es la \textbf{exponencial matricial}:
  \[
    \mathbf{P}(t) = e^{t\generator} = \sum_{k=0}^{\infty} \frac{(t\generator)^k}{k!}.
  \]
\end{teorema}

\begin{ejemplo}{Soluci\'on expl\'icita para 2 estados}
  Con la Q-matrix anterior ($\lambda, \mu > 0$), los autovalores de $\generator$ son
  $0$ y $-(\lambda+\mu)$. La soluci\'on es:
  \[
    \mathbf{P}(t) = \frac{1}{\lambda+\mu}
    \begin{pmatrix} \mu + \lambda e^{-(\lambda+\mu)t} & \lambda(1 - e^{-(\lambda+\mu)t}) \\
                    \mu(1 - e^{-(\lambda+\mu)t})    & \lambda + \mu e^{-(\lambda+\mu)t} \end{pmatrix}.
  \]
  Cuando $t \to \infty$: $\mathbf{P}(t) \to \frac{1}{\lambda+\mu}\begin{pmatrix}\mu & \lambda \\ \mu & \lambda\end{pmatrix}$,
  que es la distribuci\'on estacionaria $\pi_0 = \mu/(\lambda+\mu)$, $\pi_1 = \lambda/(\lambda+\mu)$.
\end{ejemplo}

\subsecheader{6.5 \textperiodcentered\ Distribuci\'on estacionaria de una CMTC}

\begin{definicion}{Distribuci\'on estacionaria (tiempo continuo)}
  Un vector $\boldsymbol{\pi}$ es distribuci\'on estacionaria de una CMTC si:
  \[
    \boldsymbol{\pi}\,\generator = \mathbf{0}, \quad
    \sum_{j \in \Sset} \pi_j = 1, \quad \pi_j \ge 0.
  \]
  En forma escalar: $\sum_{i \neq j} \pi_i\,q_{ij} = \pi_j\,q_j$ para todo $j$
  (flujo de entrada = flujo de salida en cada estado).
\end{definicion}

\begin{nota}
  El balance detallado en tiempo continuo es $\pi_i\,q_{ij} = \pi_j\,q_{ji}$.
  Si se satisface, $\boldsymbol{\pi}$ es estacionaria y la cadena es reversible.
\end{nota}

\subsecheader{6.6 \textperiodcentered\ Cadena embebida}

\begin{definicion}{Cadena de saltos (cadena embebida)}
  Dada una CMTC $\{X(t)\}$, la \textbf{cadena de saltos} (o \emph{cadena embebida}) es el
  proceso $\{Y_n\}_{n\ge0}$ que registra \emph{solo el estado} en el momento del $n$-\'esimo
  salto, sin tomar en cuenta el tiempo transcurrido.\\[3pt]
  La cadena de saltos es una CMTD con probabilidades de transici\'on:
  \[
    \hat{p}_{ij} = \frac{q_{ij}}{q_i} \quad (i \neq j), \qquad \hat{p}_{ii} = 0.
  \]
  As\'i, una CMTC puede describirse como: tiempos de espera Exp$(q_i)$ en cada estado $i$,
  seguidos de saltos seg\'un $\hat{p}_{ij}$.
\end{definicion}

\begin{cuidado}
  La distribuci\'on estacionaria $\boldsymbol{\pi}$ de la CMTC y la de la cadena embebida
  $\hat{\boldsymbol{\pi}}$ \emph{no coinciden} en general. La distribuci\'on estacionaria de
  la CMTC pondera adem\'as el tiempo de permanencia en cada estado:
  $\pi_i \propto \hat{\pi}_i / q_i$.
\end{cuidado}

% =============================================================================
\secheader{7}{Proceso de Nacimiento y Muerte}

\subsecheader{7.1 \textperiodcentered\ Definici\'on}

\begin{definicion}{Proceso de nacimiento y muerte}
  Un proceso de nacimiento y muerte es una CMTC con espacio de estados
  $\Sset = \{0, 1, 2, \ldots\}$ (finito o infinito numerable) donde solo se permiten
  saltos de un estado al inmediatamente adyacente:
  \begin{itemize}
    \item \textbf{Tasa de nacimiento}: $\lambda_i = q_{i,i+1} \ge 0$ (salto de $i$ a $i+1$).
    \item \textbf{Tasa de muerte}: $\mu_i = q_{i,i-1} \ge 0$ (salto de $i$ a $i-1$), con $\mu_0 = 0$.
    \item $q_i = \lambda_i + \mu_i$ (tasa total de salida del estado $i$).
  \end{itemize}
  La Q-matrix tiene la estructura tridiagonal:
  \[
    \generator =
    \begin{pmatrix}
      -\lambda_0 & \lambda_0 & 0 & 0 & \cdots \\
      \mu_1 & -(\lambda_1+\mu_1) & \lambda_1 & 0 & \cdots \\
      0 & \mu_2 & -(\lambda_2+\mu_2) & \lambda_2 & \cdots \\
      \vdots & & & \ddots &
    \end{pmatrix}.
  \]
\end{definicion}

\begin{center}
\begin{tikzpicture}[>=Latex, node distance=2.3cm,
    state/.style={circle, draw=sectionbg, fill=thmbg, line width=1pt,
                  minimum size=0.85cm, font=\small\bfseries}]
  \node[state] (0) {$0$};
  \node[state, right of=0] (1) {$1$};
  \node[state, right of=1] (2) {$2$};
  \node[state, right of=2] (3) {$3$};
  \node[right of=3, font=\large] (dots) {$\cdots$};
  \draw[->, thick, color=sectionbg]
    (0) edge[bend left=30] node[above,font=\small]{$\lambda_0$} (1)
    (1) edge[bend left=30] node[above,font=\small]{$\lambda_1$} (2)
    (2) edge[bend left=30] node[above,font=\small]{$\lambda_2$} (3)
    (3) edge[bend left=30] node[above,font=\small]{$\lambda_3$} (dots);
  \draw[->, thick, color=defborder]
    (1) edge[bend left=30] node[below,font=\small]{$\mu_1$} (0)
    (2) edge[bend left=30] node[below,font=\small]{$\mu_2$} (1)
    (3) edge[bend left=30] node[below,font=\small]{$\mu_3$} (2);
  \node[font=\small\itshape, color=sectionbg!70!black] at (4.5,-1.6)
    {Proceso de nacimiento y muerte: solo saltos al estado adyacente};
\end{tikzpicture}
\end{center}

\subsecheader{7.2 \textperiodcentered\ Balance detallado y distribuci\'on estacionaria}

\begin{teorema}{Distribuci\'on estacionaria del proceso de nacimiento y muerte}
  En un proceso de nacimiento y muerte, el balance detallado simplifica a:
  \[
    \pi_i\,\lambda_i = \pi_{i+1}\,\mu_{i+1}, \quad i = 0, 1, 2, \ldots
  \]
  (flujo entre estados $i$ e $i+1$ se equilibra en ambas direcciones).
  Iterando la recursi\'on desde $i=0$:
  \[
    \pi_n = \pi_0 \prod_{k=0}^{n-1} \frac{\lambda_k}{\mu_{k+1}}, \quad n \ge 1.
  \]
  La constante $\pi_0$ se determina por normalizaci\'on:
  $\pi_0 = \left(1 + \sum_{n=1}^{\infty} \prod_{k=0}^{n-1}\frac{\lambda_k}{\mu_{k+1}}\right)^{-1}$.
  La distribuci\'on estacionaria existe si y solo si la suma converge.
\end{teorema}

\begin{demostracion}
  Las ecuaciones de balance global para el estado $j$ son
  $\lambda_{j-1}\pi_{j-1} + \mu_{j+1}\pi_{j+1} = (\lambda_j + \mu_j)\pi_j$.
  En un proceso de nacimiento y muerte, se puede verificar que el balance detallado
  $\pi_i\lambda_i = \pi_{i+1}\mu_{i+1}$ es una soluci\'on (por inducci\'on en $i$).
  Despejando: $\pi_{i+1} = (\lambda_i/\mu_{i+1})\pi_i$, y aplicando recursivamente se
  obtiene la f\'ormula producto.
\end{demostracion}

\subsecheader{7.3 \textperiodcentered\ Proceso de nacimiento puro: conexi\'on con Poisson}

\begin{proposicion}{Proceso de Poisson como nacimiento puro}
  Si $\mu_i = 0$ para todo $i$ (no hay muertes) y $\lambda_i = \lambda$ constante
  (tasa de nacimiento homog\'enea), el proceso de nacimiento y muerte se reduce al
  \textbf{Proceso de Poisson} con intensidad $\lambda$.\\[3pt]
  En este caso no existe distribuci\'on estacionaria (el proceso crece indefinidamente):
  $\sum_{n=0}^\infty (\lambda/0)^n$ diverge.\\[3pt]
  Las probabilidades $p_n(t) = \Prob{X(t)=n \mid X(0)=0}$ satisfacen las ecuaciones
  de Kolmogorov: $p_n'(t) = -\lambda p_n(t) + \lambda p_{n-1}(t)$, con soluci\'on
  $p_n(t) = e^{-\lambda t}(\lambda t)^n/n!$ (distribuci\'on Poisson).
\end{proposicion}

\subsecheader{7.4 \textperiodcentered\ Cola M/M/1}

\begin{definicion}{Cola M/M/1}
  Una \textbf{cola M/M/1} modela un sistema de espera con:
  \begin{itemize}
    \item Llegadas: proceso de Poisson con tasa $\lambda$ (interarribos Exp$(\lambda)$).
    \item Servicio: tiempo de servicio Exp$(\mu)$.
    \item Un solo servidor, capacidad ilimitada, disciplina FIFO.
  \end{itemize}
  El n\'umero de clientes $X(t)$ es un proceso de nacimiento y muerte con
  $\lambda_i = \lambda$ y $\mu_i = \mu$ para $i \ge 1$, $\mu_0 = 0$.\\[3pt]
  El \textbf{factor de carga} es $\rho = \lambda/\mu$.
\end{definicion}

\begin{teorema}{Distribuci\'on estacionaria de M/M/1}
  La cola M/M/1 tiene distribuci\'on estacionaria si y solo si $\rho < 1$ (la tasa de
  llegadas es menor que la de servicio). Cuando se cumple:
  \[
    \pi_n = (1-\rho)\,\rho^n, \quad n = 0, 1, 2, \ldots
  \]
  Distribución geom\'etrica con par\'ametro $1-\rho$. Las medidas de rendimiento son:
  \begin{align*}
    \E{L}   &= \frac{\rho}{1-\rho} \quad \text{(n\'umero medio en el sistema)} \\
    \E{L_q} &= \frac{\rho^2}{1-\rho} \quad \text{(n\'umero medio en cola)} \\
    \E{W}   &= \frac{1}{\mu - \lambda} \quad \text{(tiempo medio en el sistema)} \\
    \E{W_q} &= \frac{\rho}{\mu - \lambda} \quad \text{(tiempo medio en cola)}
  \end{align*}
  donde se han usado las \textbf{leyes de Little}: $\E{L} = \lambda\,\E{W}$ y
  $\E{L_q} = \lambda\,\E{W_q}$.
\end{teorema}

\begin{demostracion}
  Aplicando la f\'ormula del proceso de nacimiento y muerte con $\lambda_i=\lambda$,
  $\mu_i=\mu$:
  \[
    \pi_n = \pi_0 \prod_{k=0}^{n-1}\frac{\lambda}{\mu} = \pi_0\,\rho^n.
  \]
  Normalizando: $\pi_0\sum_{n=0}^\infty \rho^n = 1 \Rightarrow \pi_0 = 1-\rho$ (requiere $\rho < 1$).\\
  El n\'umero medio: $\E{L} = \sum_{n=0}^\infty n(1-\rho)\rho^n = (1-\rho)\cdot\frac{\rho}{(1-\rho)^2}
  = \frac{\rho}{1-\rho}$.
\end{demostracion}

\begin{ejemplo}{Cola de un cajero autom\'atico}
  Clientes llegan a un ATM seg\'un un Proceso de Poisson con $\lambda = 4$ clientes/hora.
  El tiempo de servicio es Exp$(\mu=6)$ clientes/hora. Factor de carga: $\rho = 4/6 = 2/3 < 1$.
  \begin{itemize}
    \item Probabilidad de que el ATM est\'e vac\'io: $\pi_0 = 1 - 2/3 = 1/3 \approx 33\%$.
    \item N\'umero medio de clientes en el sistema: $\E{L} = (2/3)/(1/3) = 2$.
    \item Tiempo medio total en el sistema: $\E{W} = 1/(6-4) = 0.5$ horas $= 30$ min.
    \item Tiempo medio esperando en cola: $\E{W_q} = (2/3)/(6-4) = 1/3$ hora $= 20$ min.
  \end{itemize}
\end{ejemplo}

\begin{cuidado}
  Si $\rho \ge 1$ la cola \emph{no tiene r\'egimen estacionario}: la l\'inea de espera crece
  indefinidamente y el sistema se satura. La condici\'on $\rho < 1$ es indispensable para
  que exista distribuci\'on estacionaria.
\end{cuidado}

\begin{nota}
  \textbf{Leyes de Little (intuitivas):} En un sistema estable con tasa de entrada $\lambda$:
  \begin{itemize}
    \item $\E{L} = \lambda\,\E{W}$: el n\'umero promedio en el sistema es la tasa de llegadas
          por el tiempo promedio que cada cliente pasa en el sistema.
    \item $\E{L_q} = \lambda\,\E{W_q}$: an\'alogamente para la cola.
  \end{itemize}
  Las leyes de Little son universales: no dependen de la distribuci\'on de llegadas ni
  de servicio, solo de la estabilidad del sistema.
\end{nota}

% ─────────────────────────────────────────────────────────────────────────────
% Conexiones clave del módulo
% ─────────────────────────────────────────────────────────────────────────────
\vspace{10pt}
\begin{tcolorbox}[
    colback=sectionbg!8, colframe=sectionbg,
    boxrule=0.8pt, arc=4pt,
    left=8pt, right=8pt, top=6pt, bottom=6pt,
    fonttitle=\bfseries\sffamily,
    title={Conexiones clave del m\'odulo},
  ]
  \small
  \begin{itemize}
    \item Proceso de Poisson $\Rightarrow$ CMTC de nacimiento puro con $\lambda_i = \lambda$ constante
    \item Balance detallado $\pi_i q_{ij} = \pi_j q_{ji}$ $\Rightarrow$ reversibilidad $\Rightarrow$ $\pi$ f\'acil de calcular
    \item Proceso de nacimiento y muerte: $\pi_n = \pi_0\prod_{k=0}^{n-1}(\lambda_k/\mu_{k+1})$
    \item CMTC $\leftrightarrow$ cadena embebida (saltos) + tiempos de espera Exp$(q_i)$
    \item Recurrencia $\Leftrightarrow$ $\sum_n p_{ii}^{(n)} = \infty$
          \hfill (criterio de recurrencia)
    \item Teorema erg\'odico: $p_{ij}^{(n)} \to \pi_j$ para cadena erg\'odica finita
    \item $\pi_j = 1/\muu{j}$ \hfill (la distribuci\'on estacionaria es el rec\'iproco del tiempo medio de retorno)
    \item Cola M/M/1 estable $\Leftrightarrow$ $\rho = \lambda/\mu < 1$, con $\pi_n = (1-\rho)\rho^n$
  \end{itemize}
\end{tcolorbox}

\end{document}
